%=====================================================================
%  NXTCart --- IEEEtran conference manuscript (standalone)
%
%  COMPILATION STATUS: *** NOT VERIFIED ***
%  No TeX distribution (pdflatex/xelatex/lualatex/latexmk) was present on
%  the host where this file was produced, so it has NOT been compiled.
%  Treat it as unproven source. To build:
%
%      pdflatex NXTCart_IEEEtran.tex
%      pdflatex NXTCart_IEEEtran.tex          % twice for cross-references
%
%  Requires: IEEEtran.cls (from the IEEE Author Center template package)
%  Packages used, all standard in TeX Live / MiKTeX:
%      cite graphicx amsmath amssymb algorithm algpseudocode
%      booktabs array multirow url xcolor textcomp
%
%  The bibliography is inlined as thebibliography, so NO bibtex run is
%  needed. An equivalent nxtcart_cited.bib is shipped alongside for
%  reference-manager import.
%
%  Figures: each \includegraphics is commented out and replaced by a
%  visible placeholder box, so the file compiles with no image assets.
%  Uncomment the \includegraphics lines once the figures are exported
%  per NXTCart_Figures_Tables_Equations.md.
%
%  Tables I and IX are abridged here to keep the two-column layout
%  tractable; the full versions are in NXTCart_Manuscript.md.
%=====================================================================
\documentclass[conference]{IEEEtran}
\IEEEoverridecommandlockouts

\usepackage{cite}
\usepackage{graphicx}
\usepackage{amsmath,amssymb}
\usepackage{algorithm}
\usepackage{algpseudocode}
\usepackage{booktabs}
\usepackage{array}
\usepackage{multirow}
\usepackage{xcolor}
\usepackage{textcomp}
\usepackage{url}

\def\BibTeX{{\rm B\kern-.05em{\sc i\kern-.025em b}\kern-.08em T\kern-.1667em\lower.7ex\hbox{E}\kern-.125emX}}

% Placeholder figure box so the document compiles without image assets.
\newcommand{\figplaceholder}[2]{%
  \fbox{\parbox[c][#2][c]{0.95\columnwidth}{\centering\footnotesize\itshape #1}}}
\newcommand{\figplaceholderwide}[2]{%
  \fbox{\parbox[c][#2][c]{0.97\textwidth}{\centering\footnotesize\itshape #1}}}

% Pending-measurement marker: grep-able in the compiled PDF.
\newcommand{\tbm}[1]{\textcolor{red}{\texttt{[TO BE MEASURED: #1]}}}
\newcommand{\tbr}[1]{\textcolor{red}{\texttt{[TO BE RECORDED: #1]}}}

\begin{document}

\title{Settlement-Triggered, Foreground-Only Verification of\\
Scanned Products in a Barcode-First Smart Shopping Cart:\\
Prototype Architecture and Evaluation Protocol}

\author{
\IEEEauthorblockN{[Author 1 Name]\IEEEauthorrefmark{1},
                  [Author 2 Name]\IEEEauthorrefmark{1},
                  [Author 3 Name]\IEEEauthorrefmark{1},
                  [Author 4 Name]\IEEEauthorrefmark{2}}
\IEEEauthorblockA{\IEEEauthorrefmark{1}[Department], [Institution], [City], [Country]\\
Email: \{[email1], [email2], [email3]\}@[domain]}
\IEEEauthorblockA{\IEEEauthorrefmark{2}[Department], [Institution], [City], [Country]\\
Email: [email4]@[domain]}
}

\maketitle

\begin{abstract}
Self-checkout and smart-cart systems move product identification from a trained cashier to the
shopper, which creates a verification problem: the store must establish that the physical object
placed in the basket is the object whose barcode was presented. Combining barcode, weight and image
evidence is long-established prior art, so this paper does not claim that combination as a
contribution. Instead we specify and document a prototype, NXTCart, whose distinguishing design
choices are mechanisms rather than modalities: (i) visual inference is neither continuous nor
scan-triggered but is launched retrospectively from a bounded, source-timestamped frame history when
the cart's load cell reports a \emph{settled} mass change, so the evidence window is selected after
the physical event rather than guessed before it; (ii) only connected foreground regions obtained by
differencing against a \emph{pre-entry} baseline frame are embedded, so no full scene --- and
therefore no basket background --- ever reaches the recognition network; (iii) each candidate region
is scored against \emph{every} registered SKU rather than only the scanned one, and a sufficiently
confident alternate SKU anywhere in the burst \emph{vetoes} a matching region elsewhere, so an early
correct view cannot outvote a later substituted product; (iv) every verdict is bound to a transaction
and attempt identifier, and the weight and image gates remain \emph{independent and conjunctive}
rather than fused into a single score, so a visual match can never clear a mass failure. We give the
complete runtime contract: a Raspberry~Pi owns both USB cameras, barcode decoding, the catalogue, the
durable per-unit ledger and checkout permission; an Arduino/HX711 performs calibrated acquisition, a
six-sample stability test and an independent transaction-specific expected-weight comparison; a PC
runs foreground segmentation and local MobileNetV2 ONNX embedding comparison over ROS~2. We
contribute a quantitative feasibility analysis showing that, with the shipped firmware constants, the
settlement timestamp necessarily lags mechanical impact by at least approximately $0.5$~s, so the
configured window observes the product \emph{at rest} rather than in descent, and that the
configuration is self-consistent only while the impact-to-settlement delay stays below $0.90$~s ---
beyond which the selected baseline becomes contaminated by the product itself. We provide a
reproducible evaluation protocol with baselines, ablations of each mechanism, an adversarial and
nuisance scenario matrix, metric denominators and leakage controls. Results are explicitly pending.
\end{abstract}

\begin{IEEEkeywords}
smart shopping cart, product verification, barcode switching, weight reconciliation, foreground
segmentation, embedding similarity, event-triggered sensing, ROS~2, selective classification, retail
loss prevention
\end{IEEEkeywords}

%---------------------------------------------------------------------
\section*{Document Status}
This is a system-design and prototype manuscript. Every quantitative claim about recognition
accuracy, latency, power, cost or theft coverage is marked \tbm{\ldots} because the hardware
acceptance campaign of Section~\ref{sec:eval} has not been executed. The only measurements reported
as completed are the host-independent software checks of Section~\ref{sec:results}, which the authors
executed at the stated repository revision. Nothing here implies venue acceptance, peer review, store
deployment or completed field validation.

%---------------------------------------------------------------------
\section{Introduction}
\label{sec:intro}

\subsection{Problem}
Automated retail identification has been studied extensively from the perception side: surveys of
computer-vision approaches to in-store product identification~\cite{R49} and of deep learning for
retail product recognition~\cite{R50} document a large literature on detecting and classifying
merchandise. A smaller but more directly relevant line of work treats the task as \emph{verification}
rather than \emph{recognition} --- deciding whether a presented object is consistent with a declared
product identity, explicitly motivated by barcode switching and missed scans at
self-checkout~\cite{R51}. Commercial smart-cart products describe camera and weight checks on the
cart itself~\cite{R81,R82}; these are vendor descriptions, cited only as deployment context and not
as evidence of accuracy or loss reduction.

The gap addressed here is not ``can a cart see products''. It is narrower and more mechanical:
\emph{given that a shopper has scanned a barcode, which pixels constitute admissible evidence about
the object that just entered the basket, and what decision rule over those pixels is defensible when
the evidence may be incomplete?}

\subsection{Scope and what is not claimed}
Carts combining an identifier reader with a load cell and a camera are at least a decade old. A 2018
barcode-first design already combined weight checks, product images and a cloud
inventory~\cite{R84}, and a 2013 wireless-sensor cart already combined a load cell with image
comparison for billing verification~\cite{R85}. A 2026 cart contributes a load-cell
weight-reconciliation algorithm with tolerance windows and drift compensation and explicitly lists
adding camera verification to disambiguate equal-weight items as \emph{future work}~\cite{R02}. A
2025 cart fuses RGB-D vision, basket weight and RFID at the decision level and already triggers full
detection only when the weight sensor reports a significant change~\cite{R01}. Consequently this
paper claims \emph{neither} multimodal cart verification, \emph{nor} weight reconciliation, \emph{nor}
weight-triggered visual inference as novel.

The visual method is likewise deliberately conventional. Decoupling localisation from classification
and classifying by nearest-reference similarity in an embedding space --- so that a new SKU requires
only new reference photographs and no retraining --- is established retail practice~\cite{R07}, and
class-verification formulations for self-checkout products have been compared directly~\cite{R52}. We
use an unmodified pretrained MobileNetV2~\cite{R31} as a frozen feature extractor and claim no
representation-learning contribution.

\subsection{Research question}
\emph{Does transaction-bound agreement between a scanned barcode, an independently checked settled
mass change, and a foreground-only visual verification of the scanned SKU reduce accepted
substitution and billing errors relative to barcode-only, barcode-plus-weight and
barcode-plus-vision baselines, at an acceptable honest-shopper retry rate and end-to-end latency?}
This is a framing, not a result.

\subsection{Contributions}
\begin{enumerate}
\item A fully specified, source-verified runtime contract for a three-host barcode-first verification
cart on ROS~2~\cite{R54}: topic graph, message schemas, quality-of-service settings, bounded buffers,
timeout and cancellation semantics, durable per-unit state and the checkout predicate.
\emph{Evidence: static inspection of the launched source at a recorded revision.}
\item \textbf{Settlement-triggered retrospective evidence selection}, admitting a burst only if it is
dense ($\geq 12$ frames) and gap-free ($\leq 80$~ms, both edges included), and abstaining otherwise.
\emph{Implemented and unit-tested; a design choice, not a demonstrated improvement.}
\item \textbf{Foreground-only inference}: full frames are never embedded; global scene changes and
scene-sized boxes are refused. \emph{Implemented and unit-tested.}
\item \textbf{All-SKU competition with a burst-level substitution veto}. \emph{Implemented and
unit-tested; effectiveness against real lookalikes is} \tbm{substitution detection rate}.
\item \textbf{Independent, non-fused, conjunctive gates with transaction/attempt binding}, in contrast
to decision-level fusion~\cite{R01}. \emph{Implemented and unit-tested; motivated by
selective-classification reasoning~\cite{R61} rather than by measurement.}
\item A \textbf{quantitative settlement-timing feasibility analysis} (Section~\ref{sec:timing})
derived from the shipped firmware constants. \emph{Analytical; the delay distribution is}
\tbm{$\Delta_s$ distribution}.
\item A \textbf{reproducible evaluation protocol} with baselines, mechanism ablations, an adversarial
and nuisance matrix, explicit denominators and leakage controls.
\end{enumerate}
Contributions 2--5 are \emph{design choices and hypotheses}. We have not established that any
improves a measurable outcome, and Section~\ref{sec:related} states precisely which elements have
close prior art.

%---------------------------------------------------------------------
\section{Related Work}
\label{sec:related}

\subsection{Identification-centric carts}
The dominant historical architecture reads an identifier and bills it: RFID carts with mobile
billing~\cite{R10} and RFID/ZigBee automatic-billing carts~\cite{R18} establish that wireless cart
billing long predates camera-based verification. One supplied comparator uses an EM-18 reader with a
ZigBee/ESP uplink and an LCD running total~\cite{R11}; its experiments concern reader behaviour versus
distance and item removal, not SKU-level visual identity. We note a specification conflict within it:
its method text names a $125$~kHz EM-18 reader while its component table lists an $860$--$960$~MHz
band. We therefore cite it for architecture and workflow only, relying on component datasheets for
electrical facts.

A second comparator combines YOLO-based image counting with an RFID \emph{wallet} for payment and a
weight sensor comparing an aggregate basket total against an expected total~\cite{R12}. Two
distinctions matter. An RFID payment card is not a per-product tag; conflating them overstates
identification coverage. And that paper's own conclusion states that the implemented project
classifies and prices items \emph{without considering their brands} --- a material difference from
verifying a branded SKU. Earlier Springer cart work~\cite{R83} and a secure-shopping architecture with
an explicit communication protocol~\cite{R06} complete the family. None establishes that identifier
reading alone resolves substitution: a correct read is consistent with a different physical object
entering the basket.

\subsection{Weight-centric carts and reconciliation}
Load-cell carts and racks~\cite{R47}, anti-theft carts combining RFID/ZigBee with weight~\cite{R48}
and theft-control prototypes~\cite{R17} form the weight-centric family. The most directly relevant
recent work is a 2026 load-cell-assisted weight-reconciliation algorithm~\cite{R02}, which addresses
exactly the problems our mass gate faces: vibration noise on a moving cart, scan-to-settle latency and
drift. Using an HX711 at $80$~SPS with a $10$-sample moving-average filter and a tolerance window
with a constant drift offset, it reports weight-measurement error averaging $0.87\,\%$ over four
grocery items at $n=30$ trials each; a fraud-detection confusion matrix over $100$ simulated
transaction cycles (46~TP, 4~FN, 2~FP, 48~TN) giving $94.0\,\%$ accuracy, $95.8\,\%$ precision and
$92.0\,\%$ recall; an approximately $420$~ms scan-to-verification delay; and a false-positive rate
below $4\,\%$ with filtering versus $28\,\%$ without. Those are that paper's measurements under its
conditions, reproduced here only as attributed prior art.

Three observations follow. First, the headline ``$96\,\%$ weight detection accuracy'' in its abstract
is not directly derivable from its Table~I (average error $0.87\,\%$ implies a much higher figure) or
its Table~III ($94.0\,\%$); comparisons should quote the table values. Second, several of its
references could not be resolved to identifiable venues, so we cite it for its own described method
and results and do not follow its citation graph. Third, and decisively for positioning, its future
work proposes adding computer vision to resolve ambiguities among identical-weight items. NXTCart
implements that direction, making it a \emph{complement} to~\cite{R02} rather than a priority claim,
and meaning its visual gate must be evaluated specifically on the equal-mass cases where a weight
gate is provably uninformative.

\subsection{Camera-plus-weight carts}
A 2019 billing cart pairs a YOLO-style detector with an HX711 load cell for produce, reporting
3~classes and $70\,\%$ detection accuracy with the cell calibrated to $5$~kg~\cite{R04}. A 2024 cart
uses a Raspberry~Pi~4, an Arduino Mega~2560, four $100$~kg load cells via HX711 and an Edge
Impulse--trained model over 5~items; its reported evaluation is a six-position weight-distribution
test and qualitative per-item outcomes~\cite{R05}. That system introduces an ``Unknown item'' state
when the camera fails but the scale detects an addition, blocking checkout until removal --- the same
\emph{function} as NXTCart's unreconciled-mass block, which we therefore do not claim as new.

TAV-Cart~\cite{R03} is the closest \emph{design} comparator: barcode scanners, a strain-gauge load
cell and a camera-based scanner stated to ``ensure the product being placed in the cart is the right
one, thus no item switching will occur'', with cloud validation, an LED status strip, a manual
item-validation mode and UPI payment. It is explicitly an initial concept and contains no quantitative
evaluation. This is instructive: the asserted elimination of item switching is exactly the claim this
paper refuses to make, because a single camera cannot certify a fully occluded exchange
(Section~\ref{sec:limits}).

A 2026 multimodal cart~\cite{R26} uses a Raspberry~Pi~5, HX711 load cells, an Intel RealSense D435i
depth camera, YOLOv11 and A* navigation with a Firestore backend; its anti-theft logic holds checkout
on a depth intrusion without a preceding barcode trigger. Its results section is an interface
walkthrough with qualitative statements; no accuracy, latency or detection-rate figures are reported.
Cloud/edge cart baselines~\cite{R09,R24} and an AI-driven cart with tracking and inventory
forecasting~\cite{R21} complete the family; the latter is relevant only as context for \emph{proposed}
extensions, since NXTCart implements no forecasting or navigation.

The 2025 fusion cart~\cite{R01} is the closest \emph{method} comparator: RGB-D camera, load cell and
RFID reader; YOLOv8-s triggered on a stabilised weight change; a Gaussian-like weight-compatibility
score and a binary RFID support score fused into a weighted sum whose argmax is the predicted SKU,
with a low-confidence branch asking the shopper to confirm. On a self-collected set of 1000 products
($\approx$75{,}000 labelled frames, 80{,}000 event instances) it reports $98.7\,\%$ event-level top-1,
$0.8\,\%$ misidentification and $0.3$~s per item, against $93.5\,\%$ (vision-only) and $95.0\,\%$
(RFID-only), with a V/V+W/V+R/V+W+R ablation.

Four differences define NXTCart against~\cite{R01}. \emph{Objective}: \cite{R01} recognises which SKU
an event corresponds to; NXTCart verifies a scanned expectation and may refuse. \emph{Fusion}:
\cite{R01} combines modalities into one score, so visual confidence can compensate for an
inconsistent mass; NXTCart keeps the gates independent and conjunctive precisely to forbid that.
\emph{Evidence extent}: \cite{R01} selects key frames; NXTCart retains every delivered frame in the
window and refuses when the window is sparse or gapped. \emph{Region content}: \cite{R01} embeds
detector outputs; NXTCart embeds only pre-entry-baseline differences, a stronger restriction with its
own failure modes.

We record that~\cite{R01} has internal citation inconsistencies: its introduction attributes shelf
monitoring to reference [6] while its list assigns [6] to a general machine-learning paper, attributes
a multimodal e-commerce dataset to [9] while its list assigns [9] to yoghurt-cup recognition, prints
two entries numbered [7], and cites a [15] absent from a list ending at [14]. Separately, a 70/10/20
split described as performed ``at the SKU level to avoid leakage'' is difficult to reconcile with
evaluating a fixed-class YOLOv8 classifier, whose output layer cannot predict held-out classes. We
treat these as reproducibility questions, not refutations, and do not copy that paper's citation
mapping.

\subsection{Retail visual recognition and verification}
Automatic-checkout research on conveyor video follows a detect-track-count paradigm; DeepACO combines
Scaled-YOLOv4 with SORT and reports $F_1 = 0.4783$ (second place) on the 2022 AI City Challenge
Track~4 test set~A over 116 product identities, trained from 116{,}500 synthetic
composites~\cite{R27}. That $F_1$ calibrates expectations: fine-grained retail identification from
video is hard even for purpose-built pipelines. Benchmarks such as RPC~\cite{R42}, densely packed
shelf detection~\cite{R38}, hierarchical grocery datasets~\cite{R39} and digital-twin training
data~\cite{R63} define the recognition task; multi-camera kiosk classification~\cite{R08} and
context-aware fine-grained classification~\cite{R30} address viewpoint and lookalike confusion.
General-purpose representations~\cite{R58,R59} and recent SKU-level retrieval with open
vision-language models~\cite{R62} are plausible replacements for our frozen MobileNetV2 and are listed
as future comparators, not implemented alternatives.

The most informative comparator for our decision rule is~\cite{R07}: RT-DETR localisation followed by
metric-embedding similarity against a reference database, requiring no retraining for new SKUs. It
reports $79.67\,\%$ Precision@1 and $49.79\,\%$ mAP at 10 classes, degrading to $67.53\,\%$ and
$38.07\,\%$ at 40 classes, with hyperbolic embeddings outperforming Euclidean at
$24.65$--$81.09$~ms per inference and a forgetting rate under $10\,\%$, attributing the dominant error
source to classification among visually similar products rather than to localisation. Two
consequences are direct. First, since NXTCart uses the same family of mechanism (crop $\rightarrow$
frozen embedding $\rightarrow$ nearest reference) with a weaker backbone, we must not expect better
open-set behaviour --- which is itself the argument for verifying against a scanned expectation and
refusing when uncertain. Second, because visually similar products dominate the error, the lookalike
and equal-mass cases of Section~\ref{sec:eval} are the decisive experiments.

Three tasks must therefore be separated, as the literature often blurs them:
\emph{object-category recognition} (``this is a tomato''~\cite{R04}); \emph{fine-grained branded-SKU
identification} (``brand $X$, not brand $Y$''~\cite{R07,R27,R42,R30}); and \emph{verification against
a scanned expectation} (``is this consistent with the SKU just scanned, and may I
refuse?''~\cite{R51,R52}). NXTCart performs the third and uses all-SKU scoring only as a
\emph{contrastive device} within it, not as an identification claim.

\subsection{Multimodal retail sensing and fusion theory}
Shelf-side systems combining vision, weight and location --- FAIM~\cite{R45} and AIM3S~\cite{R46} ---
are the strongest open prior art on how mass and image evidence interact, including failure cases;
shelf mass is partitioned per product position, whereas a cart's single aggregate cell cannot
attribute simultaneous placements. Classical~\cite{R37} and modern~\cite{R53} fusion treatments frame
evidence combination and uncertainty, but general fusion theory does not establish that this cart's
channels are independent: the Pi's delta computation and the Arduino's comparison share one load
cell, so their agreement is \emph{computational} redundancy, not sensor redundancy. Work on
calibration~\cite{R60} and selective classification~\cite{R61} underpins our refusal behaviour: a
cosine similarity is not a probability of correct verification, and a system that may abstain should
be characterised by risk against coverage rather than by accuracy alone.

\subsection{Systems and methods background}
ROS~2~\cite{R54} supplies the distributed substrate, and its quality-of-service model~\cite{R72} is
what makes frame loss an explicit, observable design parameter rather than a hidden failure.
Detector~\cite{R34} and tracker~\cite{R35,R36} foundations are cited because the cart literature is
built on them and because NXTCart's retired predecessor used tracking; the current runtime contains no
detector and no tracker, so we report no tracking metrics.

\begin{table*}[!t]
\caption{Comparison of cart and checkout systems on verification-relevant dimensions (abridged; full
version in the supplementary material). ``FT'' = full text read, ``MD'' = metadata/abstract level
only. ``Not reported'' means absent from the source we read, not an assertion that the property is
absent. Figures are each source's own, under its own conditions, and are \emph{not} comparable across
rows.}
\label{tab:related}
\centering
\scriptsize
\renewcommand{\arraystretch}{1.15}
\begin{tabular}{@{}p{0.030\textwidth}p{0.018\textwidth}p{0.016\textwidth}p{0.115\textwidth}p{0.105\textwidth}p{0.105\textwidth}p{0.115\textwidth}p{0.085\textwidth}p{0.145\textwidth}p{0.115\textwidth}@{}}
\toprule
Ref & Yr & Rd & Identification & Vision role & Weight role & Temporal/event assoc. & Compute & Measured task \& result & Limitations \\
\midrule
\cite{R01} & 2025 & FT & RGB-D + cell + RFID, decision fusion & Primary recogniser (YOLOv8-s), fused & Gaussian compatibility vs nominal mass & Weight trigger; short window; IoU tracking & Embedded SoC, INT8 & 1000 products, $\sim$80k events: 98.7\% acc., 0.8\% misID, 0.3 s/item; V 93.5, V+W 96.1, V+R 97.3 & Occlusion, tag conflict; citation inconsistencies; SKU split vs fixed head \\
\cite{R02} & 2026 & FT & Barcode + cell only & None & Primary: delta matching, tolerance + drift offset & Scan or weight event in window $T$; 10-sample MAF & ESP32 + tablet, cloud DB & 0.87\% avg error (4 items, $n{=}30$); 100 cycles: 94.0\% acc., 95.8\% prec., 92.0\% rec.; $\sim$420 ms; FPR $<$4\% (vs 28\%) & $<$20 g exempt; equal-weight unresolved (vision = future work) \\
\cite{R03} & 2024 & FT & Barcode (laser; optional camera) & Asserted anti-switching; not evaluated & Continuous, cloud cross-check & Continuous; manual/auto validation mode & Cart CPU + cloud & None --- explicitly an initial concept & No evaluation; horizontal barcodes unscannable \\
\cite{R26} & 2026 & FT & Barcode + cell + active stereo & Depth intrusion; YOLOv11 audit & Primary: mass vs declared & Intrusion \emph{without} barcode $\rightarrow$ hold & Pi 5 + Firestore & Qualitative walkthrough; no numbers reported & No quantitative validation \\
\cite{R07} & 2025 & FT & Vision only: RT-DETR + metric embeddings & Whole pipeline; nearest reference, no retraining & None & None (shelf stills) & Workstation & 79.67\% P@1 / 49.79\% mAP @10 cls; 67.53 / 38.07 @40; 24.65--81.09 ms & Errors dominated by similar products; IoU $\approx$0.68 \\
\cite{R27} & 2022 & FT & Vision only: YOLOv4 + SORT & Whole pipeline & None & Track IDs; count inside ROI & Not on-cart & AI City 2022 T4-A: $F_1{=}0.4783$, 116 IDs & Occlusion, blur, similarity; synth-to-real gap \\
\cite{R11} & 2022 & FT & RFID (EM-18) & None & Display only & None & Nano + ESP8266 + ZigBee & Reader behaviour vs distance/removal & No SKU identity check; 125 kHz vs 860--960 MHz conflict \\
\cite{R12} & 2022 & FT & RFID on product \emph{groups} + RFID wallet & YOLOv3 counting & Aggregate total vs expected & None & On-cart + Wi-Fi & Example detections; counting & Own conclusion: prices \emph{without} brands \\
\cite{R04} & 2019 & FT & Vision only (produce) & Category classifier & Per-item mass pricing & None & Arduino + host & 3 classes; 70\% detection; cell to 5 kg & Category level only \\
\cite{R05} & 2024 & FT & Barcode + vision (Edge Impulse) & Category recogniser; ``Unknown item'' blocks checkout & 4$\times$100 kg cells; unscanned-addition detection & None & Pi 4 + Mega 2560 & 5 items; 6-position test; qualitative & Cucumber problematic; no aggregate metrics \\
\cite{R45} & 2020 & MD & Shelf vision + weight + location & One fused channel & One fused channel, per position & Shelf events & Store infra. & Inventory accuracy (abstract level) & Shelf, not cart; full text pending \\
\cite{R46} & 2019 & MD & Multimodal shelf sensing & One fused channel & One fused channel & Store events & Store infra. & Inventory monitoring (abstract level) & Shelf, not cart; full text pending \\
\cite{R51} & 2022 & MD & Declared vs observed, OCR + classification & Verification with conformal abstention & None & Per presentation & Not reported & Product verification w/ Mondrian conformal prediction & Full text pending; OCR-dependent \\
\midrule
\textbf{Ours} & 2026 & --- & \textbf{Barcode primary; vision verifies the scanned SKU, never identifies alone} & \textbf{Foreground-only ROIs vs all SKUs; abstains; vetoes on confident alternate} & \textbf{Independent conjunctive gate: Pi delta + firmware check, same cell} & \textbf{Retrospective burst from a 4 s/120-frame ring around settled mass; density+gap admissibility; txn/attempt binding} & \textbf{Pi: sensors, catalogue, ledger, checkout. Arduino: acquisition + verdict. PC: segmentation + ONNX} & \textbf{Software checks only (Sec.~\ref{sec:results}); accuracy, latency, detection} \tbm{all} & \textbf{No proof against occluded substitution; one cell cannot attribute concurrent placements; UPI QR is a request, not settlement} \\
\bottomrule
\end{tabular}
\end{table*}

%---------------------------------------------------------------------
\section{Problem Formulation and Design Requirements}
\label{sec:problem}

\subsection{Entities}
Let $\mathcal{K}$ be the set of registered SKUs. The catalogue maps a barcode to
$(\text{barcode},\text{name},\text{sku},w_e,\mathrm{tol},\text{price})$ where $w_e>0$ is the
\emph{gross packaged} mass in grams, $0<\mathrm{tol}<w_e$ and $\text{price}>0$; these constraints are
enforced at import and at every lookup. A \emph{unit} is one physical instance of a SKU in one cart
session, identified by a transaction identifier $\tau$ and an \emph{attempt} counter
$a\in\mathbb{N}^{+}$. The pair $(\tau,a)$ is the association key for every piece of evidence.

\subsection{Events and evidence}
For the $(\tau,a)$ of interest: $t_{\text{scan}}$ is a decode accepted by the scan gate; $t_s$ is the
host-clock timestamp of the first \emph{stable, calibrated} sample whose deviation from the current
anchor exceeds the noise threshold and which the ledger assigns to this unit --- \textbf{a
stable-sample timestamp, not a mechanical impact interrupt}; and
$\mathcal{F}=\{(u_i,s_i,I_i)\}$ is the PC frame ring, with $s_i$ the source (Pi-clock) timestamp
applied immediately after camera read, $u_i$ the PC monotonic receive time and $I_i$ a BGR image.

\subsection{The verification task}
Given $(\tau,a)$, the scanned SKU $k^*\in\mathcal{K}$, the settled mass delta $\Delta$ and
$\mathcal{F}$, produce a mass verdict $V_w\in\{\textsc{pass},\textsc{fail},\textsc{pending},
\textsc{checking}\}$ and an image verdict $V_c\in\{\textsc{match},\textsc{mismatch},\textsc{retry},
\textsc{unavailable}\}$ with a reason code, such that a unit becomes billable only if both gates
reach \textsc{pass} for the \emph{current} attempt, and such that a system that cannot obtain
admissible evidence \emph{abstains} rather than guessing in either direction. This is a
selective-classification formulation~\cite{R61}: coverage is a free variable and must be reported
alongside error.

\subsection{Threat and nuisance model}
\textbf{Adversarial} (must not be accepted): A1 barcode switching; A2 equal-mass substitution (mass
gate provably uninformative); A3 lookalike substitution; A4 unscanned insertion; A5 duplicate-scan
exploitation; A6 removal fraud; A7 occlusion attacks; A8 timing attacks during a camera outage,
settlement wait or restart window.

\textbf{Honest nuisances} (must not be falsely rejected, or must be recoverable): N1 hands in frame;
N2 shadows and lighting change; N3 motion blur; N4 partial occlusion by existing contents; N5 dropped
frames; N6 variable impact-to-settlement delay; N7 scale drift; N8 cart movement; N9 cancelled or
retried placements; N10 service restarts.

\subsection{Derived requirements}
Table~\ref{tab:req} lists the requirements these threats imply, together with the implementation
status of each at the revision inspected in Section~\ref{sec:impl}.

\begin{table}[!t]
\caption{Derived design requirements and their implementation status.}
\label{tab:req}
\centering
\footnotesize
\renewcommand{\arraystretch}{1.2}
\begin{tabular}{@{}p{0.025\columnwidth}p{0.36\columnwidth}p{0.10\columnwidth}p{0.39\columnwidth}@{}}
\toprule
ID & Requirement & Rationale & Status in current code \\
\midrule
R1 & Evidence associable to exactly one $(\tau,a)$ & A5, A6, A8 & Implemented: $(\tau,a)$ keys on both hosts; completed-key cache; identity mismatch $\rightarrow$ \textsc{retry} \\
R2 & Mass and image decisions must not compensate & A2, A3 & Implemented: conjunctive gates; a camera \textsc{match} cannot clear a mass \textsc{fail} \\
R3 & Evidence window chosen \emph{after} the physical event & N6 & Implemented: retrospective extraction from a bounded ring \\
R4 & Incomplete visual evidence abstains & N5, N6 & Implemented: density and gap tests; missing baseline $\rightarrow$ \textsc{retry} \\
R5 & Background pixels must not enter the network & N2, N4 & Implemented: foreground-only ROIs; global change refused \\
R6 & Early correct view must not outvote a later substitution & A2, A3, A7 & Implemented: all-SKU scoring + burst-level veto \\
R7 & Unexplained mass blocks checkout until accounted mass returns & A4 & Implemented: unassigned-weight issue latches the checkout reason \\
R8 & Loss of observation invalidates trust & A8, N10 & Implemented: camera-generation change, heartbeat gap and scale restart force reconciliation \\
R9 & Restart must not resurrect unfinished units as verified & A8, N10 & Implemented: SQLite WAL journal; non-\textsc{verified} units failed on load \\
R10 & Concurrency limits must be physical & A5, N8 & Implemented: overlapping unsettled placements flagged ambiguous, never guessed \\
R11 & Payment authorisation $\neq$ payment request & --- & \emph{Partial}: revision-gated UPI QR; \textbf{no provider confirmation exists} \\
R12 & Light curtains, multiple cells, IMU, navigation, forecasting & future & \textbf{Not implemented}; documented as planned design only \\
\bottomrule
\end{tabular}
\end{table}

%---------------------------------------------------------------------
\section{System Architecture}
\label{sec:arch}

\subsection{Host responsibilities}
Fig.~\ref{fig:arch} shows three hosts with a deliberately rigid split, enforced mechanically by an
import-boundary check in the test runner. The \textbf{Arduino + HX711} performs calibrated
acquisition at ${\sim}10$~Hz, a six-sample stability test and an \emph{independent}
transaction-specific expected-mass comparison, with calibration persisted in EEPROM and no auto-tare
at boot. The \textbf{Raspberry~Pi} owns both USB cameras, barcode/QR decoding, the read-only
catalogue, the serial bridge, the durable per-unit ledger, the touchscreen and \textbf{checkout
permission}; it loads no model, embedding database, reference photograph or credential. The
\textbf{PC} owns the frame ring, foreground segmentation, the MobileNetV2 ONNX database and the
verdict publisher; it imports nothing from the Pi package, decides no prices and holds no cart state.

\textbf{The PC is advisory; the Pi is authoritative.} This is the most important architectural
property: a compromised, lagging or restarted verification service can withhold approval but cannot
grant it, alter a bill or authorise payment.

\begin{figure}[!t]
\centering
% \includegraphics[width=\columnwidth]{fig1_architecture}
\figplaceholder{\textbf{Fig.~1 placeholder.} Export the Mermaid source in the supplementary
artifacts file (\texttt{mmdc -i fig1.mmd -o fig1.pdf}) and uncomment the \texttt{includegraphics}
line above. Three swimlanes: Arduino/HX711 $\mid$ Raspberry~Pi (authority) $\mid$ PC (advisory).}{46mm}
\caption{NXTCart host boundaries and topic graph. Dashed arrows carry advisory verdicts. The single
\texttt{/cart/placement\_settled} edge is the only inference trigger. The import boundary (no
\texttt{ml}/\texttt{torch}/\texttt{onnxruntime}/\texttt{nxtcart\_vision} in the Pi package, no
\texttt{smart\_cart} in the PC package) is enforced by a static check in CI.}
\label{fig:arch}
\end{figure}

\subsection{Topic graph and message contract}
Table~\ref{tab:topics} gives the contract. A deliberate asymmetry: image transport is
\emph{best-effort} on both ends, so frames may be lost by design~\cite{R72}. Rather than hiding that
loss, the burst admissibility test of Section~\ref{sec:burst} makes it an explicit, reported outcome.

\begin{table*}[!t]
\caption{ROS~2 interface contract. Message definitions are byte-identical in both workspaces, checked
by CI.}
\label{tab:topics}
\centering
\footnotesize
\renewcommand{\arraystretch}{1.15}
\begin{tabular}{@{}p{0.175\textwidth}p{0.115\textwidth}p{0.145\textwidth}p{0.175\textwidth}p{0.305\textwidth}@{}}
\toprule
Topic & Type & Publisher $\rightarrow$ Subscriber & QoS & Role \\
\midrule
\texttt{/cart/camera/frame} & \texttt{sensor\_msgs/Image} (bgr8) & Pi camera $\rightarrow$ PC & Pub: sensor data (best effort); Sub: best effort, depth 30 & Source-stamped frames; \texttt{frame\_id = verification\_camera/<uuid>} carries the camera generation \\
\texttt{/cart/scanned\_item} & \texttt{ScannedItem} & Pi scanner $\rightarrow$ authority, PC & Reliable, depth 50 & Barcode + \texttt{scan\_id}. \textbf{Logged by the PC; never starts inference} \\
\texttt{/cart/verification\_request} & \texttt{VerificationRequest} & authority $\rightarrow$ PC & Reliable, depth 50 & \texttt{action} $\in$ \{add, remove, cancel\}; creates or cancels context \\
\texttt{/cart/placement\_settled} & \texttt{VerificationRequest} & authority $\rightarrow$ PC & Reliable, depth 50 & \texttt{started\_at} = settled-sample stamp. \textbf{The only inference trigger} \\
\texttt{/cart/vision\_event} & \texttt{VisionEvent} & PC $\rightarrow$ authority & Reliable, depth 50 & Verdict + reason code + scores, bound to $(\tau,a,\text{action})$ \\
\texttt{/cart/vision\_health} & \texttt{String} (JSON) & PC $\rightarrow$ authority & Reliable, depth 10 & 2~Hz heartbeat: \texttt{ready}, \texttt{generation}, \texttt{inflight}, \texttt{collecting}, \texttt{buffered\_frames}, \texttt{camera\_age\_seconds} \\
\texttt{/cart/scene\_anomaly} & \texttt{String} & PC $\rightarrow$ authority & Reliable, depth 10 & Camera restart notice $\rightarrow$ forces reconciliation \\
\texttt{/cart/weight\_sample} & \texttt{WeightSample} & bridge $\rightarrow$ authority & Reliable, depth 50 & \texttt{stream\_id}, \texttt{seq}, \texttt{grams}, \texttt{stable}, \texttt{calibrated}, receive stamp \\
\texttt{/cart/weight\_expectation} & \texttt{WeightExpectation} & authority $\rightarrow$ bridge $\rightarrow$ Arduino & Reliable, depth 50 & \texttt{E,txn,baseline\_g,expected\_g,tolerance\_g} \\
\texttt{/cart/weight\_verdict} & \texttt{WeightVerdict} & bridge $\rightarrow$ authority & Reliable, depth 50 & Firmware's independent \textsc{pass}/\textsc{fail} with measured delta \\
\texttt{/cart/ui\_display} & \texttt{CartState} & authority $\rightarrow$ GUI & Reliable, transient-local, depth 1 & Per-unit statuses, totals, \texttt{revision}, \texttt{checkout\_allowed}, \texttt{checkout\_reason} \\
\texttt{/cart/\{remove,checkout\}\_request}, \texttt{/cart/checkout\_quote}, \texttt{/cart/alerts} & \texttt{String} & GUI $\leftrightarrow$ authority & Reliable, depth 10 & Selected removal; revision-gated checkout quote \\
\bottomrule
\end{tabular}
\end{table*}

\subsection{Durable per-unit state}
Each unit is persisted to SQLite in WAL mode with \texttt{synchronous=FULL}, alongside an append-only
audit table and a monotonically increasing \texttt{revision}. The derived display status is
\textsc{flagged} if any of $\{V_w,V_c,V_{\mathrm{ard}}\}$ is \textsc{fail}, \textsc{verified} if
$V_w=V_c=\textsc{pass}$, and \textsc{verifying} otherwise, with a separate \textsc{removing} branch.
On process start, any unit that is not \textsc{verified} is forced to \textsc{fail} --- a restart can
never resurrect an unfinished unit as billable (R9). See Fig.~\ref{fig:state}.

\begin{figure}[!t]
\centering
% \includegraphics[width=\columnwidth]{fig4_state}
\figplaceholder{\textbf{Fig.~4 placeholder.} Per-unit state machine; export from the Mermaid
\texttt{stateDiagram-v2} source in the supplementary artifacts file.}{42mm}
\caption{Per-unit verification state. The two gates are independent and conjunctive;
\textsc{verified} requires both \textsc{pass} for the \emph{current} attempt, and \textsc{fail} on
either gate is terminal for that attempt. The firmware's \textsc{checking} intermediate exists so a
Pi-side mass pass cannot become \textsc{pass} before the independent firmware comparison arrives.
Recovery from \textsc{flagged} is always via selected removal, which increments the attempt counter.}
\label{fig:state}
\end{figure}

\subsection{Checkout predicate}
Checkout is allowed only when all of the following hold, evaluated in this order so the shopper gets
the most actionable reason: (1)~the cart is non-empty; (2)~the scale is ready and not in recovery;
(3)~a fresh ($<2.5$~s) stable calibrated sample exists; (4)~the PC heartbeat is fresh ($<2.5$~s);
(5)~no global issue is latched; (6)~every unit is \textsc{verified}; (7)~$|g_{\text{latest}} -
\sum_u \mathrm{measured\_g}(u)| \le 10$~g; (8)~a merchant UPI address is configured; (9)~the quote's
revision equals the ledger's. Condition~(7) is a \emph{second}, aggregate mass check layered on the
per-unit checks: a basket of correctly verified units can still be refused if the total does not
reconcile. Conditions (8)--(9), plus the GUI's own $1.5$~s freshness and \texttt{request\_id} checks,
mean a displayed quote closes as soon as the cart changes.

\subsection{Shopper experience versus physical concurrency}
Accepted scans appear in the bill immediately with per-unit Weight/Camera indicators; both gates run
in the background, and a flagged unit does \emph{not} lock the scanner. This separates two properties
that are often conflated: \textbf{GUI responsiveness} (rows appear at once and are never blocked by a
pending or failed earlier unit) from \textbf{physical concurrency restrictions}, which are a hardware
observability limit enforced regardless of interface latency. With one aggregate load cell the system
cannot attribute mass to individual simultaneous placements, so the ledger enforces: more than one
unit awaiting weight $\Rightarrow$ \emph{all} pending units flagged ambiguous and a global overlap
issue latched; a new scan during a removal $\Rightarrow$ immediate mass-gate failure; a scan while the
scale is not ready or not fresh $\Rightarrow$ immediate mass-gate failure; and a mass step not
attributable to exactly one pending unit $\Rightarrow$ an ``unassigned weight'' issue with the anchor
advanced so the evidence cannot be reused by a later scan. A second placement may legitimately begin
as soon as the previous unit's mass has settled, even while its image inference is still running; up
to eight verifications may be in flight.

%---------------------------------------------------------------------
\section{Methodology and Algorithms}
\label{sec:method}

\subsection{Frame ingestion and the bounded ring}
Each accepted frame is stored as $(u_i,s_i,I_i)$ in a deque bounded to $N_{\mathrm{buf}}=120$ entries
and pruned to $T_{\mathrm{buf}}=4.0$~s of receive time. A frame is rejected if its encoding, dtype or
dimensions are invalid, if its source-clock age falls outside $[-0.5,5.0]$~s, or if its source stamp
does not strictly exceed the previous stamp for the same camera generation (a replayed physical frame
must not supply a second sample). A change of \texttt{frame\_id} is a \emph{camera generation change}:
the ring is cleared, all in-flight work returns \textsc{retry}, and a scene-anomaly notice forces
reconciliation. Memory and bandwidth follow from the configuration:
\begin{align}
M_{\mathrm{ring}} &= N_{\mathrm{buf}} W H \cdot 3 = 120\cdot640\cdot480\cdot3\ \mathrm{B} = 110.6\ \mathrm{MB}\\
B_{\mathrm{net}}  &= W H \cdot 3 f = 640\cdot480\cdot3\cdot30\ \mathrm{B/s} \approx 221\ \mathrm{Mb/s}
\end{align}
Both are \emph{analytical} figures for raw BGR payload only; DDS, IP and framing overhead, worker
snapshots and ROS internal buffers add to both. Achieved figures: \tbm{achieved publish rate, receive
rate, wire bandwidth, PC RSS}.

\subsection{Settlement-triggered burst assembly}
\label{sec:burst}
On \texttt{/cart/placement\_settled} with $(\tau,a)$, the node validates that the settlement matches
the requested transaction (scan id, barcode, resolved SKU and action must all agree, else
\textsc{retry}), that the SKU resolves in the embedding database (else \textsc{unavailable}), that the
settlement stamp age lies in $[-0.1,T_{\mathrm{buf}}]$ (else \textsc{retry}, flagging clock skew), and
that the in-flight limit is not exceeded. It then registers a \emph{collection} rather than
snapshotting immediately, because the post-event frames have not necessarily arrived and the two
topics may be reordered. With $t_s$ the settlement stamp,
\begin{align}
\mathcal{W}(t_s) &= [\,t_s-\delta_{\mathrm{pre}},\ t_s+\delta_{\mathrm{post}}\,],\quad
  \delta_{\mathrm{pre}}{=}0.40\,\mathrm{s},\ \delta_{\mathrm{post}}{=}0.15\,\mathrm{s}\\
b(t_s) &= t_s - \delta_{\mathrm{base}},\qquad \delta_{\mathrm{base}} = 0.90\,\mathrm{s}\\
\mathcal{B} &= \{(u,s,I)\in\mathcal{F} : b-\delta_{\mathrm{tol}} \le s \le b\},\
  \delta_{\mathrm{tol}}{=}0.15\,\mathrm{s}
\end{align}
Collection completes when the newest received source stamp passes $t_s+\delta_{\mathrm{post}}$, or
when a transport grace period $T_{\mathrm{grace}}=0.75$~s expires. The burst
$\mathcal{S}=\{(u,s,I)\in\mathcal{F} : t_s-\delta_{\mathrm{pre}} \le s \le t_s+\delta_{\mathrm{post}}\}$
is \emph{admissible} only if
\begin{align}
\mathcal{B} &\neq \emptyset\\
|\mathcal{S}| &\ge n_{\min} = 12\\
\max_i\big(\tilde{s}_{i+1}-\tilde{s}_i\big) &\le g_{\max} = 80\,\mathrm{ms}
\end{align}
where $\tilde{s}=[t_s-\delta_{\mathrm{pre}},s_1,\dots,s_{|\mathcal{S}|},t_s+\delta_{\mathrm{post}}]$.
Including both window edges is deliberate: a burst dense in the middle but missing its first or last
$100$~ms would otherwise pass. At 30~FPS an admissible burst is normally 16--17 frames. The chosen
baseline is $\max(\mathcal{B})$ by source stamp, i.e.\ the newest frame at or before $b(t_s)$.
Algorithm~\ref{alg:burst} gives the procedure.

\begin{algorithm}[!t]
\caption{Settlement-triggered burst assembly (PC)}
\label{alg:burst}
\footnotesize
\begin{algorithmic}[1]
\Statex \textbf{On \texttt{/cart/placement\_settled}} with $m=(\tau,a,\ldots,t_s)$:
\If{$(\tau,a)$ completed, running or collecting} \State \Return \Comment{idempotent} \EndIf
\State $c \gets \textsc{PendingPop}(\tau,a)$ \textbf{or} $\textsc{ContextFrom}(m)$
\If{$\textsc{Identity}(c)\neq\textsc{Identity}(m)$}
  \State \textsc{Publish}(\textsc{retry}, ``settlement mismatch''); \Return \EndIf
\If{$\textsc{Resolve}(\mathrm{sku},\mathrm{barcode})=\bot$}
  \State \textsc{Publish}(\textsc{unavailable}, ``unknown SKU''); \Return \EndIf
\State $\mathrm{age}\gets\textsc{HostClockNow}()-t_s$
\If{$\mathrm{age}\notin[-0.1,T_{\mathrm{buf}}]$}
  \State \textsc{Publish}(\textsc{retry}, ``stale settlement; check clock sync''); \Return \EndIf
\If{$|\mathrm{collecting}|+|\mathrm{running}|\ge 8$}
  \State \textsc{Publish}(\textsc{retry}, ``queue full''); \Return \EndIf
\State $w\gets\max(0,\delta_{\mathrm{post}}-\mathrm{age})$
\State $\mathrm{collecting}[(\tau,a)]\gets(c,t_s,\mathrm{now}{+}w,\mathrm{now}{+}w{+}T_{\mathrm{grace}})$
\Statex \textbf{On every 20\,Hz tick:}
\ForAll{$(\tau,a)\in\mathrm{collecting}$}
  \State $\mathrm{covered}\gets(\text{newest stamp}\ge t_s+\delta_{\mathrm{post}})$
  \If{$\mathrm{now}<\mathrm{ready}$ \textbf{or} ($\neg\mathrm{covered}$ \textbf{and}
       $\mathrm{now}<\mathrm{deadline}$)}
    \State \textbf{continue} \Comment{keep receiving; never snapshot on the callback} \EndIf
  \State remove $(\tau,a)$ from $\mathrm{collecting}$
  \State $\mathcal{B}\gets\{f : t_s{-}\delta_{\mathrm{base}}{-}\delta_{\mathrm{tol}}\le
         \mathrm{stamp}(f)\le t_s{-}\delta_{\mathrm{base}}\}$
  \State $\mathcal{S}\gets\{f : t_s{-}\delta_{\mathrm{pre}}\le\mathrm{stamp}(f)\le
         t_s{+}\delta_{\mathrm{post}}\}$
  \If{$\mathcal{B}=\emptyset$}
    \State \textsc{Publish}(\textsc{retry}, ``no pre-entry baseline''); \textbf{continue} \EndIf
  \State $\mathrm{gaps}\gets\textsc{Diffs}([t_s{-}\delta_{\mathrm{pre}}]\|
         \mathrm{stamps}(\mathcal{S})\|[t_s{+}\delta_{\mathrm{post}}])$
  \If{$|\mathcal{S}|<n_{\min}$ \textbf{or} $\max(\mathrm{gaps})>g_{\max}$}
    \State \textsc{Publish}(\textsc{retry}, ``incomplete burst''); \textbf{continue} \EndIf
  \State \textsc{SubmitToSingleWorker}(\textsc{VerifyBurst}, $\textsc{Newest}(\mathcal{B}),
         \mathcal{S}, c$)
\EndFor
\end{algorithmic}
\end{algorithm}

\subsection{Foreground extraction}
For a baseline $B$ and a burst frame $I$ of identical geometry, the mask is the thresholded
\emph{per-channel maximum} absolute difference --- chosen over a grayscale difference so a change
between two similarly bright but differently coloured surfaces is preserved:
\begin{align}
D(\mathbf{x}) &= \max_{c\in\{B,G,R\}}\big|B_c(\mathbf{x})-I_c(\mathbf{x})\big|\\
M(\mathbf{x}) &= \mathbb{1}\!\left[D(\mathbf{x})>\tau_d\right],\qquad \tau_d = 25\\
M' &= \mathrm{close}_{5\times5}\!\big(\mathrm{open}_{3\times3}(M)\big)
\end{align}
using standard binary morphology~\cite{R66}. Two refusals follow, both preventing background from
reaching the network:
\begin{align}
\textstyle\sum_{\mathbf{x}} M'(\mathbf{x}) &> \rho H W,\ \rho=0.65 \Rightarrow \text{refuse}\\
(x_1{-}x_0)(y_1{-}y_0) &> \rho H W \Rightarrow \text{refuse}
\end{align}
where the padded box uses $p=0.12$:
\begin{equation}
\begin{aligned}
x_0 &= \max(0,x{-}\lceil pw\rceil), & y_0 &= \max(0,y{-}\lceil ph\rceil),\\
x_1 &= \min(W,x{+}w{+}\lceil pw\rceil), & y_1 &= \min(H,y{+}h{+}\lceil ph\rceil).
\end{aligned}
\end{equation}
Components below $A_{\min}=800$~px are discarded as noise; if more than eight components survive the
frame is refused as ``too many changed regions'' --- many disjoint changes indicate camera shake rather
than one inserted product, and failing closed is preferred to silently discarding a possibly
substituted object. Crucially, components are \emph{never merged}: a hand and a product at opposite
corners stay two tight crops rather than becoming one scene-sized rectangle (Fig.~\ref{fig:fg}).
Algorithm~\ref{alg:fg} states the procedure.

For a \textbf{selected removal} the roles invert: the departing product is present in the baseline and
absent from the burst frames, so the difference is computed with the burst frame as reference and the
\emph{crops are taken from the baseline image}. Cropping the now-empty basket would embed background
and reintroduce exactly the leakage this design removes.

\begin{figure}[!t]
\centering
% \includegraphics[width=\columnwidth]{fig5_foreground}
\figplaceholder{\textbf{Fig.~5 placeholder.} Foreground-only ROI pipeline; export from the Mermaid
flowchart in the supplementary artifacts file.}{40mm}
\caption{Foreground-only ROI pipeline. Three distinct refusal paths (global scene change, oversized
box, too many regions) replace what would otherwise be a full-frame fallback. Zero regions across the
whole burst yields \textsc{mismatch}/\texttt{NO\_ITEM\_SEEN}, which is tested \emph{before}
\texttt{SCENE\_CHANGED}.}
\label{fig:fg}
\end{figure}

\begin{algorithm}[!t]
\caption{Foreground ROI extraction (no full frame is ever embedded)}
\label{alg:fg}
\footnotesize
\begin{algorithmic}[1]
\Require baseline $B$, frame $I$, $\tau_d,A_{\min},p,\rho$,
         $\mathrm{action}\in\{\textsf{placed},\textsf{removed}\}$
\If{$\mathrm{action}=\textsf{removed}$} \State $\textsc{Swap}(B,I)$
  \Comment{crop from the baseline, never the empty basket} \EndIf
\If{$\mathrm{shape}(B)\neq\mathrm{shape}(I)$}
  \State \Return \textsc{reject}(``frame geometry changed'') \EndIf
\State $M\gets\mathbb{1}[\max_c|B_c-I_c|>\tau_d]$
\State $M\gets\textsc{Close}_{5\times5}(\textsc{Open}_{3\times3}(M))$
\If{$\textsc{NonZero}(M)>\rho H W$}
  \State \Return \textsc{reject}(``global scene change'') \EndIf
\State $\mathcal{C}\gets$ external contours of $M$, area descending
\ForAll{$\kappa\in\mathcal{C}$}
  \If{$\textsc{Area}(\kappa)<A_{\min}$} \State \textbf{continue} \EndIf
  \State $\beta\gets\textsc{Clip}(\textsc{Pad}(\textsc{BoundingRect}(\kappa),p))$
  \If{$\textsc{Area}(\beta)>\rho H W$}
    \State \Return \textsc{reject}(``box covers too much of the scene'') \EndIf
  \State append $\beta$, $\textsc{Crop}(I,\beta)$ \Comment{components are NEVER merged}
\EndFor
\If{$|\mathrm{crops}|>8$}
  \State \Return \textsc{reject}(``too many changed regions'') \Comment{fail closed} \EndIf
\State \Return $\mathrm{crops},\mathrm{boxes}$
\end{algorithmic}
\end{algorithm}

\subsection{Embeddings and per-region scoring}
Each crop is resized to $224\times224$, converted BGR$\rightarrow$RGB, scaled to $[0,1]$ and
normalised with the ImageNet statistics used by the torchvision ONNX export, then transposed to NCHW.
The embedding is the 1280-dimensional global-average-pooling output --- the standard
``strip the classification head'' construction --- L2-normalised:
\begin{equation}
\mathbf{z} = \mathrm{GAP}\big(\mathrm{MobileNetV2}_{\mathrm{body}}(\mathrm{norm}(\mathrm{crop}))\big),
\quad \mathbf{e} = \mathbf{z}/\lVert\mathbf{z}\rVert_2
\end{equation}
The classifier logits are \emph{not} the reference space, and the extractor resolves the pooling layer
at start-up by probing for a 1280-element output, failing startup if none exists. Reference vectors
$\mathbf{r}_{k,j}$ are built offline from product crops by the same preprocessing, so references and
live crops share one framing domain. Appearance similarity is the maximum cosine over that SKU's
references:
\begin{equation}
s(k,\mathbf{e}) = \mathrm{clip}\!\left(\max_j\langle\mathbf{r}_{k,j},\mathbf{e}\rangle,-1,1\right)
\end{equation}
Colour evidence, when present, is an 80-bin HSV fingerprint: an $18\times4$ hue$\times$saturation
histogram over \emph{chromatic} pixels ($S\ge90$, $V\ge60$) concatenated with an 8-bin brightness
histogram over \emph{achromatic} pixels ($S<90$, $V\ge40$), so white, grey and dark packaging are
distinguishable. With $\mathbf{h}(\cdot)$ the L2-normalised fingerprint,
\begin{equation}
c(k,\mathrm{crop}) = \mathrm{clip}\!\left(\max_j\langle\mathbf{h}_{k,j},
  \mathbf{h}(\mathrm{crop})\rangle,0,1\right),\ \ \text{ok} \iff c\ge0.30
\end{equation}
and, if a SKU declares an expected mean BGR colour $\boldsymbol{\mu}_k$,
\begin{equation}
c_{\mathrm{bgr}} = 1-\frac{\lVert\mathrm{mean}(\mathrm{crop})-\boldsymbol{\mu}_k\rVert_2}{255\sqrt{3}},
  \ \ \text{ok} \iff c_{\mathrm{bgr}}\ge0.30
\end{equation}
with the reported colour score the minimum of the two channels when both exist. Apparent shape
evidence uses only the crop's aspect ratio against reference ratios $\rho_j$:
\begin{equation}
\alpha = \frac{\max(h,w)}{\min(h,w)},\quad
\text{size\_ok} \iff \exists j:\max\!\left(\tfrac{\alpha}{\rho_j},\tfrac{\rho_j}{\alpha}\right)\le3.0
\label{eq:aspect}
\end{equation}
\textbf{Equation~(\ref{eq:aspect}) is explicitly not a physical size measurement.} The camera is
uncalibrated, with no depth channel and no known scale, so $\alpha$ is an image-plane ratio that
varies with distance, pose and crop tightness. It is a weak consistency filter and is treated as such.

\subsection{All-SKU competition and the burst-level veto}
For each region all $|\mathcal{K}|$ SKUs are ranked by $s(k,\mathbf{e})$ with a deterministic
tie-break on the SKU key. With $k_1$ the top-1 SKU, $s_1$ its score, $s_2$ the runner-up,
$m=s_1-s_2$ and $k^*$ the scanned SKU:
\begin{equation}
\begin{cases}
\textsc{mismatch}/\texttt{SWAP\_DETECTED} & k_1\neq k^* \wedge s_1\ge\theta_{\mathrm{swap}}{=}0.70\\
\textsc{match} & k_1 = k^* \wedge s_1\ge\theta{=}0.72 \wedge m\ge\mu{=}0.04 \wedge \text{ok}\\
\textsc{mismatch}/\texttt{LOW\_CONFIDENCE} & \text{otherwise}
\end{cases}
\end{equation}
Aggregation over the burst is a \emph{lexicographic precedence}, not a vote:
\begin{equation}
\mathcal{P} = \begin{cases}
\mathcal{S}\mathrm{w} & \mathcal{S}\mathrm{w}\neq\emptyset\\
\mathcal{M} & \mathcal{S}\mathrm{w}=\emptyset,\ \mathcal{M}\neq\emptyset\\
\mathcal{L} & \text{otherwise}
\end{cases}
\qquad \hat{r} = \arg\max_{r\in\mathcal{P}}\mathrm{score}(r)
\label{eq:veto}
\end{equation}
Hence a single sufficiently confident alternate SKU anywhere in the burst overrides every matching
region: an early, correct, hand-held view of the scanned product cannot certify a different object
that was actually released. If no region produced a valid embedding the outcome is
\textsc{unavailable}/\texttt{NO\_EMBEDDING}; if no region existed it is
\textsc{mismatch}/\texttt{NO\_ITEM\_SEEN}. Algorithm~\ref{alg:verify} states the procedure.

Note the asymmetry between $\theta_{\mathrm{swap}}=0.70$ and $\theta=0.72$: accusation is marginally
\emph{easier} than approval. That is intentional in a fail-closed design, but it also means a
near-threshold lookalike can produce a \texttt{SWAP\_DETECTED} against an honest shopper. The
honest-shopper cost is \tbm{false \texttt{SWAP\_DETECTED} rate on genuine placements}.

\begin{algorithm}[!t]
\caption{Burst verification with all-SKU competition and veto}
\label{alg:verify}
\footnotesize
\begin{algorithmic}[1]
\Require scanned SKU $k^*$; per-frame region lists; reference DB over $\mathcal{K}$;
         $\theta,\mu,\theta_{\mathrm{swap}}$
\If{$k^*\notin\mathcal{K}$}
  \State \Return \textsc{unavailable}(``no reference embeddings'') \EndIf
\State $\mathcal{E},\mathcal{S}\mathrm{w},\mathcal{M},\mathcal{L}\gets[\,],[\,],[\,],[\,]$
\For{each frame index $i$, each region index $j$}
  \State $x\gets\mathcal{R}_i[j]$
  \If{$x$ invalid \textbf{or} $\min(h,w)<8$}
    \State record \textsc{unavailable}; \textbf{continue} \EndIf
  \State $\mathbf{e}\gets\textsc{Normalise}(\textsc{Embed}(x))$
    \Comment{single worker: \texttt{cv2.dnn} never re-entered}
  \State $\mathrm{ranked}\gets\textsc{SortDesc}[(s(k,\mathbf{e}),c,\mathrm{ok},k):k\in\mathcal{K}]$
    \Comment{EVERY SKU}
  \State $(s_1,c_1,\mathrm{ok}_1,k_1)\gets\mathrm{ranked}[0]$;\ \
         $s_2\gets\mathrm{ranked}[1].s$;\ \ $m\gets s_1-s_2$
  \If{$k_1\neq k^*$ \textbf{and} $s_1\ge\theta_{\mathrm{swap}}$}
    \State append to $\mathcal{S}\mathrm{w}$
  \ElsIf{$k_1=k^*$ \textbf{and} $s_1\ge\theta$ \textbf{and} $m\ge\mu$ \textbf{and} $\mathrm{ok}_1$}
    \State append to $\mathcal{M}$
  \Else \State append to $\mathcal{L}$ \EndIf
  \State record $(i,j,k_1,s_1,s(k^*),k_2,s_2,m,c_1,\mathrm{ok}_1)$ in $\mathcal{E}$
\EndFor
\State $\mathcal{P}\gets\mathcal{S}\mathrm{w}$ \textbf{or} $\mathcal{M}$ \textbf{or} $\mathcal{L}$
  \Comment{Eq.~(\ref{eq:veto}): a confident swap vetoes every match}
\State \Return $\arg\max_{r\in\mathcal{P}}\mathrm{score}(r)$ with $\mathcal{E}$, else
       \textsc{unavailable}/\texttt{NO\_EMBEDDING}
\end{algorithmic}
\end{algorithm}

\subsection{Independent mass verification}
Two comparisons are made over the same physical signal. The firmware maintains a six-sample ring at
${\sim}10$~Hz and declares
\begin{equation}
\mathrm{stable} \iff \mathrm{calibrated} \wedge |\mathrm{ring}|{=}6 \wedge
  (\max\mathrm{ring}-\min\mathrm{ring}) \le \sigma = 4.0\,\mathrm{g}
\label{eq:stable}
\end{equation}
publishing the window mean when stable and the instantaneous reading otherwise. On a Pi-supplied
expectation it emits, once a stable reading deviates from the baseline by more than $5$~g,
\begin{equation}
\Delta_{\mathrm{ard}} = \mathrm{reading}-\mathrm{baseline\_g},\quad
  \textsc{pass} \iff |\Delta_{\mathrm{ard}}-w_e| \le \mathrm{tol}
\end{equation}
Independently, the Pi's ledger requires
\begin{equation}
\Delta_{\mathrm{pi}} = g-\mathrm{anchor},\quad
  \textsc{pass} \iff |\Delta_{\mathrm{pi}}-w_e| \le \mathrm{tol}
\end{equation}
A unit's mass gate becomes \textsc{pass} only when both hold; while the firmware verdict is
outstanding the gate sits in \textsc{checking}, and a firmware \textsc{fail} is terminal and cannot be
overwritten by a later pass. A missing verdict times out to \textsc{fail} after the 20~s ledger
deadline and the firmware's own 20~s expectation timeout. This is redundancy of \emph{computation and
of expectation plumbing}, not of sensing (Section~\ref{sec:limits}).

\subsection{Does the configured window capture the physical insertion?}
\label{sec:timing}
This is the question the configuration most depends on, and it is answerable analytically from the
shipped constants without any hardware. Let $t_0$ be mechanical impact and $t_q\ge t_0$ the moment
the reading first enters a band of width $\le\sigma=4$~g and stays there. Because
Eq.~(\ref{eq:stable}) requires six consecutive $100$~ms samples inside the band, the first sample that
can be marked stable is the sixth, i.e.\ no earlier than $t_q+0.5$~s, and in practice
$t_q+0.5\ldots0.6$~s depending on sample phase. The Pi stamps the sample on serial receipt, adding
transmission and USB-CDC latency on the order of milliseconds. Writing
$\Delta_s = t_s-t_0 = (t_q-t_0)+\Delta_{\mathrm{stab}}+\Delta_{\mathrm{serial}}$ with
$\Delta_{\mathrm{stab}}\in[0.5,0.6]$~s:
\begin{align}
\text{descent inside the window} &\iff \Delta_s \le \delta_{\mathrm{pre}} = 0.40\,\mathrm{s}
  \label{eq:c27}\\
\text{baseline strictly pre-entry} &\iff \Delta_s < \delta_{\mathrm{base}} = 0.90\,\mathrm{s}
  \label{eq:c28}
\end{align}

\textbf{Condition~(\ref{eq:c27}) is unsatisfiable with the shipped firmware}, because
$\Delta_{\mathrm{stab}}$ alone exceeds $\delta_{\mathrm{pre}}$. The configured burst therefore does
\emph{not} observe the product in flight; it observes the product \emph{at rest}, and possibly the
hand withdrawing. This is not necessarily a defect --- a settled product is still foreground with
respect to a pre-entry baseline, and the absence of the hand reduces occlusion --- but it has three
consequences that must be stated rather than assumed away.

First, the 16--17 frames of an admissible burst are \emph{near-duplicates} of one static scene.
Density buys robustness against packet loss and shutter noise, not viewpoint diversity. Claims that
depend on temporal diversity (for example, that a burst sees multiple faces of a package) are
unsupported.

Second, the margin before baseline contamination is thin. Combining~(\ref{eq:c28}) with
$\Delta_s \ge \Delta_{\mathrm{stab}} \approx 0.5\ldots0.6$~s leaves roughly \textbf{0.30--0.40~s of
mechanical decay budget}. Because the selected baseline is the newest frame at $t_s-0.90$~s, once
$\Delta_s>0.90$~s the baseline itself already contains the product at rest; the difference mask then
collapses and the system emits \textsc{mismatch}/\texttt{NO\_ITEM\_SEEN} --- a \textbf{false reject
with a misleading reason code}, not a timeout.

Third, the self-consistent operating band for the current configuration is
\begin{equation}
0.40\,\mathrm{s} < \Delta_s < 0.90\,\mathrm{s}
\label{eq:band}
\end{equation}
which must be verified empirically for the actual mechanism, product mass range and mounting
stiffness: \tbm{distribution of $\Delta_s$ over $\ge200$ drops, per mass decile, on the mounted
cart}. If the upper tail crosses $0.90$~s, the correct responses are to widen $\delta_{\mathrm{base}}$
(and $T_{\mathrm{buf}}$ accordingly), to reduce $\Delta_{\mathrm{stab}}$ by raising the HX711 sample
rate or shortening the stability window at the cost of noise immunity, or to add a \emph{separately
defined impact timestamp} --- for example a first-derivative threshold crossing published alongside the
settled sample. Lowering the confidence thresholds to compensate for missing evidence is explicitly
not an acceptable response, because it trades a false reject for a false accept. Fig.~\ref{fig:timing}
shows the regimes.

\begin{figure}[!t]
\centering
% \includegraphics[width=\columnwidth]{fig2_timing}
\figplaceholder{\textbf{Fig.~2 placeholder.} Draw in TikZ from the coordinate specification in the
supplementary artifacts file: axis $-1.05$ to $+0.90$~s relative to $t_s$; baseline search window
$[-1.05,-0.90]$; burst window $[-0.40,+0.15]$ with 17 ticks at $33.3$~ms; transport-grace deadline
$+0.90$; and the $t_{\mathrm{impact}}$ uncertainty bar spanning $0.5$--$0.9$~s left of the origin.
Shade the band of Eq.~(\ref{eq:band}).}{48mm}
\caption{Settlement-triggered burst timing at the configured defaults. $t_{\mathrm{impact}}$ lies
$0.5$--$0.9$~s before $t_s$, i.e.\ the descent happens \emph{before} the window opens. Region~I
($\Delta_s\le0.40$~s, descent captured) is unreachable with the shipped firmware; Region~II
($0.40<\Delta_s<0.90$~s) is the valid at-rest band; Region~III ($\Delta_s\ge0.90$~s) contaminates the
baseline and yields a mislabelled false reject.}
\label{fig:timing}
\end{figure}

\subsection{Outcome codes and their consequences}
Table~\ref{tab:codes} lists the implemented outcomes. Two implementation details deserve explicit
statement because they are not obvious from the documentation. First, \texttt{NO\_ITEM\_SEEN} is
tested \emph{before} \texttt{SCENE\_CHANGED}: if \emph{every} frame in a burst is refused as a global
scene change, the region count is zero and the result is \textsc{mismatch}/\texttt{NO\_ITEM\_SEEN},
not \textsc{retry}/\texttt{SCENE\_CHANGED}. Anyone interpreting the diagnostic manifests should
therefore read \texttt{NO\_ITEM\_SEEN} as ``no usable foreground region was isolated anywhere in the
burst'', a class that includes the total-illumination-change case; the per-frame \texttt{rejected}
field in the manifest distinguishes the two causes. The safety direction is identical either way ---
both paths refuse. Second, the Pi collapses the four-valued image verdict to a binary gate ---
\textsc{pass} only for \textsc{match} with \texttt{approved} true, \textsc{fail} otherwise --- so
\textsc{retry} and \textsc{unavailable} are \emph{not} retried automatically; the shopper's recovery
path is selected removal and re-scan, which creates a fresh transaction with an incremented attempt.
Both behaviours are recorded here as observed properties of the evaluated revision. No change to
either is proposed in this paper.

\begin{table}[!t]
\caption{Verdict and reason codes, as implemented, with their effect on the Pi.}
\label{tab:codes}
\centering
\footnotesize
\renewcommand{\arraystretch}{1.2}
\begin{tabular}{@{}p{0.20\columnwidth}p{0.14\columnwidth}p{0.56\columnwidth}@{}}
\toprule
Code & Verdict & Produced when / effect \\
\midrule
\texttt{MATCH} & \textsc{match} & Target top-1 above $\theta$ with margin $\ge\mu$ and colour/size
  ok, and no swap anywhere in the burst. $V_c\gets\textsc{pass}$ \\
\texttt{SWAP\_DETECTED} & \textsc{mismatch} & An alternate SKU scored $\ge\theta_{\mathrm{swap}}$;
  takes precedence over all matches. $V_c\gets\textsc{fail}$, reason carries the observed SKU \\
\texttt{LOW\_CONFIDENCE} & \textsc{mismatch} & Regions exist but none qualifies.
  $V_c\gets\textsc{fail}$ \\
\texttt{NO\_ITEM\_SEEN} & \textsc{mismatch} & \emph{No} foreground region in any burst frame; also
  reached when every frame was refused. $V_c\gets\textsc{fail}$ \\
\texttt{SCENE\_CHANGED} & \textsc{retry} & At least one frame yielded regions \emph{and} at least one
  was refused. $V_c\gets\textsc{fail}$ \\
\texttt{INFERENCE\_ERROR} & \textsc{unavailable} & The matcher raised inside the worker \\
\texttt{NO\_EMBEDDING} & \textsc{unavailable} & Regions existed but none produced a valid embedding \\
(text reason) & \textsc{retry} & Missing baseline; sparse or gapped burst; stale settlement; identity
  mismatch; queue full; camera restart; inference deadline; settlement never arrived \\
(text reason) & \textsc{unavailable} & Scanned barcode/SKU not in the embedding database \\
\bottomrule
\end{tabular}
\end{table}

\subsection{Concurrency, cancellation and liveness}
All ROS callbacks run in the default mutually-exclusive callback group; all segmentation, inference
and diagnostic disk writes run on a \emph{single} worker thread, so \texttt{cv2.dnn} is never entered
concurrently and camera/heartbeat callbacks never block behind inference. No worker publishes or
mutates node state: results are collected on the 20~Hz timer. The frame ring, the in-flight map
($\le8$) and the completed-key cache ($\le4096$, FIFO-evicted) are all bounded.

Cancellation is handled at three points, each covered by a regression test: a \texttt{cancel} action
removes pending and collecting entries and marks a running future expired; an expired future's result
is discarded when it completes; and a verdict for an already-completed key is suppressed. A camera
generation change invalidates everything in flight. Heartbeats continue during inference and tolerate
a camera age up to $5.0$~s even after old frames have been pruned, so a brief delivery stall does not
invalidate already-verified units --- but a genuine camera restart does, via
\texttt{/cart/scene\_anomaly}.

Fig.~\ref{fig:seq} summarises one successful placement end to end.

\begin{figure}[!t]
\centering
% \includegraphics[width=\columnwidth]{fig3_sequence}
\figplaceholder{\textbf{Fig.~3 placeholder.} Export the Mermaid \texttt{sequenceDiagram} from the
supplementary artifacts file. Six participants: scanner camera, authority, firmware, verification
camera, verifier, touchscreen.}{44mm}
\caption{Message sequence for one successful placement. A barcode scan creates verification context
but never starts inference; the PC logs it and waits. Inference begins only on
\texttt{/cart/placement\_settled}. The PC does not snapshot on the settlement callback, because
post-event frames may not yet have arrived and the two topics may be delivered out of order.}
\label{fig:seq}
\end{figure}

%---------------------------------------------------------------------
\section{Implementation}
\label{sec:impl}

The system was inspected and tested at repository revision
\texttt{28aa58d3ceff6a7d03a2d208d9d4065606af85f8} (11~October~2026). All statements in this section
derive from the launched source and configuration at that revision.

\subsection{Hardware and software inventory}
HX711 with DT$\rightarrow$pin~3 and SCK$\rightarrow$pin~2 at $115200$~baud~\cite{R74}; the
\texttt{bogde/HX711} Arduino library~\cite{R75}; calibration persisted as an EEPROM
record~\cite{R79} via \texttt{ZERO} then \texttt{CAL,<grams>}, each averaging ten raw reads, with no
auto-tare at boot. Stability uses $\mathrm{WINDOW}{=}6$ samples,
$\mathrm{STABLE\_SPAN\_G}{=}4.0$~g and ${\sim}100$~ms sample gating; the firmware trigger is
$\mathrm{CHANGE\_G}{=}5.0$~g. Load-cell capacity, class and excitation are deployment choices not
fixed by the code. Both hosts run Ubuntu~24.04 with ROS~2 Jazzy,
\texttt{ROS\_DOMAIN\_ID=42} and \texttt{rmw\_cyclonedds\_cpp} on a shared trusted LAN with DDS
UDP/multicast permitted. The scanner camera is a USB webcam at $640\times480$ with
\texttt{BUFFERSIZE=1} under V4L2, decoded with pyzbar/ZBar~\cite{R78}; the verification camera is a
USB webcam at $640\times480$ and $30$~FPS \emph{requested} (the driver-reported rate is logged at
open; no MJPG FourCC is requested, so raw-format USB bandwidth may limit the achieved rate):
\tbm{achieved frame rate}. Exact Pi model, RAM, PC CPU and RAM: \tbr{host specifications}.

The model is MobileNetV2~\cite{R31}, \texttt{mobilenetv2-7.onnx} from the ONNX model zoo
(torchvision export, opset~7, ${\approx}13.9$~MB), with the 1280-D GAP output resolved at start-up by
probing; preprocessing is $224\times224$, BGR$\rightarrow$RGB, $/255$, mean $(0.485,0.456,0.406)$,
std $(0.229,0.224,0.225)$, NCHW --- identical in the offline database builder and the live matcher.
Inference uses \texttt{cv2.dnn.readNetFromONNX} with \texttt{ENGINE\_CLASSIC} where available,
\texttt{DNN\_BACKEND\_OPENCV} and \texttt{DNN\_TARGET\_CPU}~\cite{R77}; there is no GPU, no ONNX
Runtime and no remote API in the live path. The colour fingerprint is the 80-bin format of
Section~\ref{sec:method} at format version~2; a stale colour format disables colour, which makes the
verifier \emph{stricter}, never laxer. The catalogue is read-only SQLite validated on every lookup,
and scanner-supplied price/name fields are \emph{not} trusted. Durable state is SQLite WAL with
\texttt{synchronous=FULL}. Payment is a locally generated, revision-gated
\texttt{upi://pay?pa\&pn\&am\&cu=INR} QR~\cite{R80} --- \textbf{a payment request, not a
provider-confirmed settlement}.

\subsection{Configured parameters}
Table~\ref{tab:params} lists the configured defaults. Internal consistency is enforced at start-up
rather than trusted: the node refuses to launch unless all durations are finite and positive, all
counts are positive integers, $\delta_{\mathrm{base}}>\delta_{\mathrm{pre}}$,
$T_{\mathrm{buf}}>\delta_{\mathrm{base}}+\delta_{\mathrm{post}}+T_{\mathrm{grace}}$,
$n_{\min}\ge2$, $N_{\mathrm{buf}}\ge n_{\min}$, $0\le p\le0.5$, $0<\rho<1$ and $\tau_d\le255$. With
the shipped values $T_{\mathrm{buf}}=4.0>1.80$ and $N_{\mathrm{buf}}=120=4.0\,\mathrm{s}\times30$~FPS,
so the ring is exactly sized for the configured window at the requested frame rate --- a 15~FPS camera
would halve ring occupancy and \emph{fail} the density test rather than silently degrade.

\begin{table}[!t]
\caption{Configured defaults. These require empirical validation; none is an optimised or
experimentally justified value.}
\label{tab:params}
\centering
\footnotesize
\renewcommand{\arraystretch}{1.1}
\begin{tabular}{@{}llll@{}}
\toprule
Group & Parameter & Value & Controls \\
\midrule
\multirow{6}{*}{Thresholds}
 & \texttt{similarity\_threshold} $\theta$ & 0.72 & min.\ top-1 cosine \\
 & \texttt{competitive\_margin} $\mu$ & 0.04 & min.\ top-1 $-$ runner-up \\
 & \texttt{swap\_threshold} $\theta_{\mathrm{swap}}$ & 0.70 & swap accusation \\
 & \texttt{colour\_threshold} & 0.30 & min.\ colour score \\
 & \texttt{size\_ratio\_tolerance} & 3.0 & apparent aspect ratio \\
 & \texttt{removal\_margin} & 0.04 & open-SKU removal margin \\
\midrule
\multirow{2}{*}{Buffer}
 & \texttt{buffer\_seconds} & 4.0 s & ring retention \\
 & \texttt{buffer\_max\_frames} & 120 & ring capacity \\
\midrule
\multirow{7}{*}{Burst}
 & \texttt{burst\_before} $\delta_{\mathrm{pre}}$ & 0.40 s & window start \\
 & \texttt{burst\_after} $\delta_{\mathrm{post}}$ & 0.15 s & window end \\
 & \texttt{baseline\_before} $\delta_{\mathrm{base}}$ & 0.90 s & baseline target \\
 & \texttt{baseline\_tolerance} $\delta_{\mathrm{tol}}$ & 0.15 s & baseline search \\
 & \texttt{burst\_min\_frames} $n_{\min}$ & 12 & density test \\
 & \texttt{burst\_max\_gap} $g_{\max}$ & 0.08 s & gap test (edges incl.) \\
 & \texttt{transport\_grace} $T_{\mathrm{grace}}$ & 0.75 s & late-delivery wait \\
\midrule
\multirow{4}{*}{Foreground}
 & \texttt{difference\_threshold} $\tau_d$ & 25 & per-channel threshold \\
 & \texttt{foreground\_min\_area} $A_{\min}$ & 800 px & min.\ contour area \\
 & \texttt{roi\_padding} $p$ & 0.12 & ROI padding \\
 & \texttt{max\_roi\_fraction} $\rho$ & 0.65 & refusal fractions \\
\midrule
\multirow{5}{*}{Scheduling}
 & \texttt{camera\_loss\_timeout} & 5.0 s & heartbeat frame age \\
 & \texttt{settlement\_timeout} & 18.0 s & pending-context wait \\
 & \texttt{result\_timeout} & 6.0 s & inference deadline \\
 & \texttt{max\_inflight} & 8 & concurrent verifications \\
 & \texttt{debug\_keep} & 100 & retained diagnostics \\
\midrule
\multirow{5}{*}{Pi ledger}
 & \texttt{verification\_timeout} & 20.0 s & per-unit deadline \\
 & \texttt{sensor\_timeout} & 2.5 s & scale/heartbeat freshness \\
 & \texttt{noise\_g} & 5.0 g & mass-step noise gate \\
 & \texttt{empty\_tolerance\_g} & 10.0 g & empty/reconciliation \\
 & \texttt{max\_units} & 200 & cart capacity \\
\midrule
Scanner & \texttt{rearm\_seconds} & 0.8 s & observed absence to re-arm \\
\midrule
\multirow{3}{*}{Firmware}
 & \texttt{WINDOW} & 6 & stability ring length \\
 & \texttt{STABLE\_SPAN\_G} $\sigma$ & 4.0 g & stability span \\
 & \texttt{CHANGE\_G} & 5.0 g & firmware trigger \\
\bottomrule
\end{tabular}
\end{table}

\subsection{Scan admission}
Repeated scans of the same code are gated by \emph{observed absence}, not by a timer: a code re-arms
only after being continuously absent from the decoder output for $0.8$~s. A camera outage calls an
explicit gap handler that resets absence timers, so an outage can never be mistaken for the code
leaving view --- otherwise unplugging the scanner camera while holding a code would re-arm it. If more
than one code is visible simultaneously, no scan is emitted and the shopper is asked to present one
code at a time.

\subsection{Diagnostic artifacts}
Each worker evaluation writes, to a UUID-named directory, the baseline image, every evaluated full
frame, every ROI and a \texttt{manifest.json} containing source timestamps, window bounds, bounding
boxes, per-region target/top-1/runner-up scores, margins, colour scores, evidence flags, verdicts and
reason codes. Retention is bounded to 100 transactions, and \emph{only} directories matching the
writer's own \texttt{nxtcart\_<UTC>\_<32-hex>} pattern are eligible for deletion; symlinked
directories are never followed and the parent is re-resolved immediately before recursive removal. A
full or unwritable disk is reported as an explicit error and \emph{cannot} change the decision --- a
regression test asserts that a \textsc{mismatch} cannot become a \textsc{match} because diagnostics
failed. These artifacts are the intended instrument for the measurement campaign and are why
Section~\ref{sec:timing} is empirically checkable: the manifest records \texttt{settled\_ns},
\texttt{window\_start\_ns}, \texttt{window\_end\_ns}, \texttt{baseline\_ns}, the realised frame count
and the realised maximum gap for every transaction. They are \textbf{diagnostic model outputs, not
labelled ground truth}, and a later cancellation or timeout can suppress publication of a verdict
whose artifacts were already written.

\subsection{Clock synchronisation assumptions}
Both camera stamps and settlement stamps originate on the \emph{Pi} clock; the PC compares them
against its \emph{own} ROS clock to compute ages and uses its monotonic clock only for scheduling and
retention. The deployment therefore requires NTP/chrony synchronisation, and simulated time must never
be mixed with wall-clock sensor stamps. Skew is not silently absorbed: a settlement stamp outside
$[-0.1,4.0]$~s returns \textsc{retry} naming clock synchronisation, and a frame outside
$[-0.5,5.0]$~s is dropped. The residual exposure is a \emph{slowly} accumulating skew smaller than
those windows, which would bias baseline and burst selection without triggering any guard:
\tbm{Pi$\leftrightarrow$PC offset and drift over a shopping session}.

\subsection{Deployment-boundary enforcement and retained code}
The test runner parses every Python source on both hosts and fails the build if the Pi package imports
\texttt{ml}, \texttt{torch}, \texttt{tensorflow}, \texttt{onnxruntime} or \texttt{nxtcart\_vision}, or
if the PC package imports \texttt{smart\_cart}; it also asserts that the two copies of
\texttt{smart\_cart\_interfaces} are byte-identical. The architectural split is thus a checked
invariant rather than a convention.

The repository deliberately retains historical material that the integrated runtime does \emph{not}
launch, and this must be stated to avoid architectural confusion. The PC package's console entry
points are exactly \texttt{vision\_verify\_node} and \texttt{capture\_references}, and
\texttt{vision\_verify\_node} imports only the matcher and burst modules. Consequently the MOG2 motion
detector, the three-window ``journey'' verifier, the geometric tracking pipeline, the zone/custody/
identity trackers, the legacy suites, the vision-API backend utilities and the JPEG-carrying
\texttt{CameraFrame} message are \emph{not part of the live path}. Three Pi modules are retired stubs
that raise on invocation. Two documentation artifacts remain stale at this revision and are superseded
by the current vision guide: the top-level README's description of ``3--5 full-frame motion views'',
and the setup guide's JPEG transport, region parameters, nine-frame scan-to-rest journey and
vision-API backend. Readers reproducing this work should resolve any architectural question from the
verification configuration, the two launch files and the three live PC modules.

%---------------------------------------------------------------------
\section{Experimental Design and Evaluation Protocol}
\label{sec:eval}

No part of this protocol has been executed. It is specified at a level intended to be directly
runnable, and every table it produces is left with explicit placeholders.

\subsection{Ground-truth units and denominators}
Three nested units of analysis, each with its own denominator, must be reported separately; reporting
only the first is the most common error in this literature. \emph{Region/frame level}: one evaluated
ROI, denominator = ROIs reaching the matcher; useful only for diagnosing the visual channel, and
\emph{not} the system's output. \emph{Transaction level}: one $(\tau,a)$ attempt, denominator =
attempts of the relevant class; the natural unit for false accept/reject. \emph{Session level}: one
complete trip ending in a checkout decision, denominator = sessions; the unit that matters
commercially, because the system's output is a bill. Ground truth is recorded by a human operator
independently of the system --- true SKU, true physical action, true mass on a reference scale,
intended scenario class and wall-clock time --- and joined to logs by $(\tau,a)$ \emph{after} the
session, never by trusting the system's own identification.

\subsection{Baselines}
All baselines run on the same physical placements and the same recorded evidence, so differences
reflect the decision rule rather than the trial set. \textbf{B0} barcode only: bill every accepted
scan. \textbf{B1} barcode + weight: require both mass checks, ignore all image evidence; reproduces
the family of~\cite{R02}. \textbf{B2} barcode + vision: require the image gate, accept any mass.
\textbf{B3} full system: conjunctive gates plus global reconciliation and the checkout predicate.
\textbf{B4} fused score (contrast): replace the conjunction with a weighted sum of a mass-consistency
score and the visual score, argmax above a threshold, in the style of~\cite{R01}. B4 is important for
honesty: the conjunctive design is a \emph{hypothesis}, and the experiment must be able to show that
fusion would have been better.

\subsection{Ablations of the claimed mechanisms}
Each ablation removes exactly one mechanism from B3, replaying the \emph{same} recorded bursts offline
so only the decision changes. \textbf{A-FG}: embed the full frame instead of foreground ROIs
(tests contribution~3). \textbf{A-ALL}: score only the scanned SKU, threshold absolutely (4).
\textbf{A-VETO}: keep all-SKU scoring but let any \textsc{match} win (4). \textbf{A-MARGIN}: set
$\mu{=}0$. \textbf{A-COLOUR}: disable colour and aspect. \textbf{A-DENSITY}: accept any burst with
$\ge1$ frame and no gap test (2). \textbf{A-BASE}: use the frame immediately before settlement as
baseline. \textbf{A-SAMPLE}: decimate the burst to five evenly spaced frames. \textbf{A-TRIG}: trigger
on the scan with a fixed forward window instead of on settlement (2). The retired full-frame and
journey-based pipelines may be evaluated \emph{only} as separately identified experiments, re-run under
the same conditions on the same recorded data; their historical development numbers are not comparable
and are not reported as results.

\subsection{Scenario matrix}
Table~\ref{tab:scen} abridges the matrix; the full 29-row version is in the supplementary material.
``In scope'' means the current code has an implemented, unit-tested handling path; it does \emph{not}
mean the handling has been validated on hardware. Each scenario requires at least 30 genuine
placements per SKU plus deliberately chosen confusable pairs, over an empty and a partly filled
basket, all relevant orientations, normal shopping speed and the real store LAN. Thirty placements per
cell is an initial bench set, \emph{not} statistical certification: with 30 trials and zero observed
failures the one-sided $95\,\%$ Clopper--Pearson upper bound on the failure rate is still
approximately $9.5\,\%$, which must be stated whenever a cell reads ``0 failures''.

\begin{table}[!t]
\caption{Scenario matrix (abridged). Out-of-scope rows are retained deliberately.}
\label{tab:scen}
\centering
\footnotesize
\renewcommand{\arraystretch}{1.15}
\begin{tabular}{@{}p{0.04\columnwidth}p{0.40\columnwidth}p{0.36\columnwidth}p{0.10\columnwidth}@{}}
\toprule
\# & Scenario & Expected behaviour & Scope \\
\midrule
S1 & Correct product, correct mass, clear view & \textsc{match} + mass \textsc{pass} & Yes \\
S2 & Correct product, wrong mass & Mass \textsc{fail}; image cannot clear it & Yes \\
S3 & Equal-mass different SKU & Mass may pass; \texttt{SWAP\_DETECTED} expected & Yes \\
S4 & Lookalike variant, same brand & Swap or low confidence; never \textsc{match} & Yes \\
S5 & Unscanned insertion & Unassigned-mass issue latches & Yes \\
S6--S8 & Duplicate scan; two scans before settling; two objects together & Ambiguous; no attribution attempted & Yes \\
S9 & Selected removal of the correct unit & Removed only after both removal gates & Yes \\
S10 & Select expensive row, remove cheap equal-mass item & Removal identity mismatch; no discount & Yes \\
S11--S15 & Hands, shadows, blur, partial occlusion, global lighting change & Refusal or low confidence; never a false \textsc{match} & Yes \\
S16 & \textbf{Fully occluded substitution} & \textbf{Cannot be resolved; must remain unverified} & \textbf{No} \\
S17--S19 & Dropped frames; camera stall 1--4 s; camera restart & \textsc{retry}; verified units preserved; mandatory reconciliation & Yes \\
S20 & Settlement delay beyond $0.90$~s & Baseline contamination $\rightarrow$ false reject & Yes (the failure \emph{is} the measurement) \\
S21 & Scale drift across a session & Anchor advances per attributed step & Partial (no drift compensation) \\
S22--S27 & Cart pushed; cancel mid-verification; duplicate/late messages; authority restart; PC stopped; Arduino unplugged & Withhold, idempotent, fail closed; no fake approval & Yes \\
S28 & Items below the mass noise floor & Mass gate cannot discriminate & \textbf{No} \\
S29 & Loose produce / variable-weight goods & Not supported by the catalogue schema & \textbf{No} \\
\bottomrule
\end{tabular}
\end{table}

\subsection{Metrics}
Table~\ref{tab:metrics} states each metric with its denominator. The risk--coverage curve is the
honest way to report a system that may abstain~\cite{R61}, and thresholds should be reported as a
curve rather than a single operating point. Because cosine similarity is not a calibrated
probability~\cite{R60}, no row should be presented as a confidence or a probability of correctness.

\begin{table}[!t]
\caption{Metrics, each with its denominator stated.}
\label{tab:metrics}
\centering
\footnotesize
\renewcommand{\arraystretch}{1.15}
\begin{tabular}{@{}p{0.34\columnwidth}p{0.30\columnwidth}p{0.28\columnwidth}@{}}
\toprule
Metric & Definition & Denominator \\
\midrule
False accept rate & Adversarial attempts reaching \textsc{verified} & Adversarial attempts of that class \\
Honest false reject rate & Honest attempts flagged on first attempt & Honest attempts \\
Retry rate & Attempts ending \textsc{retry}/\textsc{unavailable} & All attempts (reported \emph{separately} from rejects) \\
Intervention rate & Sessions needing reconciliation or staff action & Sessions \\
Exact basket / total correctness & Final billed multiset / total equals ground truth & Sessions \\
Substitution detection rate & Swap attempts yielding \texttt{SWAP\_DETECTED} with the \emph{correct} observed SKU & Swap attempts \\
False accusation rate & Honest attempts yielding \texttt{SWAP\_DETECTED} & Honest attempts \\
Verification latency (p50/p95) & $t_{\mathrm{verdict}}-t_s$ & Attempts producing a verdict \\
End-to-end latency (p50/p95) & $t_{\mathrm{verdict}}-t_{\mathrm{scan}}$ & As above (reported \emph{in addition}) \\
Achieved rate / gap distribution & From manifest realised counts and gaps & Admissible + inadmissible bursts \\
Burst admissibility rate & Bursts passing the density and gap tests & Settlement events \\
Settlement delay $\Delta_s$ & $t_s - t_{\mathrm{impact}}$ via an independent reference & Drops \\
Mass error & Measured delta $-$ reference-scale mass & Placements \\
Resource use / network & CPU, RSS, ms per ROI; bytes/s on the wire & Sampled \\
Risk--coverage curve & FAR vs fraction decided, as $(\theta,\mu,\theta_{\mathrm{swap}})$ vary & Attempts \\
\bottomrule
\end{tabular}
\end{table}

\subsection{Tuning/test separation and leakage controls}
Thresholds, colour and aspect tolerances and all burst timings are tuned on a \emph{development}
split and frozen before the test campaign; the frozen values are reported. Reference photographs,
tuning data and test data come from \emph{disjoint capture sessions}, disjoint \emph{physical units}
of each product and disjoint lighting conditions. Adjacent frames from the same burst must never be
split across tuning and test --- the unit of partition is the \emph{capture session}, never the frame.
Sixteen near-duplicate frames from one burst carry roughly one burst's worth of independent
information (Section~\ref{sec:timing}), so treating frames as independent samples would inflate every
interval. Lookalike and equal-mass pairs are selected and held out \emph{before} any tuning. Each
reported figure states the number of sessions, attempts and physical units contributing to it, with
Clopper--Pearson intervals for proportions.

\subsection{Instrumentation plan}
An independent impact reference is required to measure $\Delta_s$; a contact microphone or
accelerometer on the basket, logged on the Pi clock, is sufficient and need not be part of the
product. Ground truth for identity comes from the operator's log; the visual channel's own manifest
provides per-region diagnostics but must not be used as its own ground truth. Wire bandwidth is
measured with a packet counter on the PC interface, not inferred from~(2).

%---------------------------------------------------------------------
\section{Results and Discussion}
\label{sec:results}

\subsection{What was actually measured}
The only completed measurements are host-independent software checks, executed at revision
\texttt{28aa58d} on Windows~11 with Python~3.14.7, OpenCV~5.0.0 and NumPy~2.5.2 --- a development
host, \emph{not} the ROS Jazzy target. Source parse, package XML, interface parity and
deployment-boundary checks: \textbf{PASS} (85 Python sources parsed; both interface copies
byte-identical; no boundary violation). Integration and regression suites: \textbf{103 tests, OK}.
Reference-handling suite: \textbf{4 tests, OK}. Scripted verifier suite: \textbf{all checks passed}.

Of the 103 tests, 35 exercise the current burst pipeline directly (8 foreground-segmentation tests
using real OpenCV, 6 all-SKU/veto ranking tests, 18 node-level burst/cancellation/health tests, 3
camera-node tests), 26 exercise the Pi ledger, 11 the audit regressions, 4 the ROS adapters, 2 the
catalogue and serial contracts, 9 the matcher's database handling and 4 the Pi scan/verdict binding.
\textbf{Twelve tests target retired components} (4 journey-window, 5 geometric-merge calibration, 3
MOG2 motion) and are not evidence about the current pipeline.

\subsection{What these results do and do not establish}
They establish that the source parses, that the two hosts' message contracts are identical, that the
deployment boundary holds, and that the implemented decision logic behaves as specified on constructed
inputs --- including that an empty scene and a global lighting change never invoke the embedding
matcher, that a four-frame packet-loss gap returns \textsc{retry} rather than a decimated decision,
that a later confident swap vetoes an earlier match while every region is still evaluated, that a
cancelled or over-deadline job cannot publish an approval, and that a diagnostic disk failure cannot
upgrade a \textsc{mismatch} to a \textsc{match}.

They establish \textbf{nothing} about live DDS behaviour or wire latency (ROS callbacks are driven
through message-schema stubs and a patched clock); achieved camera frame rate; recognition accuracy on
real packaging; HX711 accuracy, linearity, creep or drift on a mounted basket; the distribution of
$\Delta_s$; native \texttt{colcon} builds; Arduino compilation; or any security property. Neither
\texttt{colcon} nor the Arduino toolchain was available on the development host, so both native builds
remain required bring-up steps rather than completed validations. The test environment's OpenCV~5.0.0
also exceeds the \texttt{>=4.10,<5} range pinned for the ROS target, so the suites were exercised
against a newer DNN stack than the deployment one.

\subsection{Discussion}
\textbf{Foreground-only inference removes a specific failure mode and introduces another.} By
construction no basket background reaches the network, so the matcher cannot score a scene against a
SKU and the ``empty basket scores against the target'' pathology that motivated the design is
structurally impossible --- a regression test pins exactly that case. The cost is that the method's
entire sensitivity now rests on the quality of one baseline frame. Baseline differencing cannot
distinguish a product from a hand holding it, cannot see a fully occluded object, and degrades under
shadows, reflections and any scene change between baseline and burst. The refusal rules convert those
degradations into refusals rather than confident errors --- the right safety direction --- but refusals
have a shopper cost that is \tbm{honest retry rate}.

\textbf{All-SKU competition is a contrastive device, not an identification claim.} Requiring the
scanned SKU to be top-1 with a margin is strictly stronger than an absolute threshold, and the margin
is what makes a lookalike pair expensive to pass. But the open-set evidence from the closest
methodological comparator is sobering: a stronger pipeline than ours reports $67.5\,\%$ Precision@1 at
only 40 classes, with errors concentrated on visually similar products~\cite{R07}. We should therefore
expect our top-1 ranking to be unreliable as a catalogue grows, and $\theta_{\mathrm{swap}}=0.70$ to
generate false accusations on honest lookalike placements. This is the sharpest open risk in the
design, it is measurable (Table~\ref{tab:metrics}), and it is why the paper does not present the veto
as a solved mechanism.

\textbf{Independence costs coverage and is the point.} Refusing to let a visual match offset a mass
failure guarantees that the attack class the mass channel \emph{can} see is never argued away by the
channel that cannot. The symmetric cost is that two independent gates multiply their false-reject
rates, so B3 will have a strictly higher honest-retry rate than either B1 or B2. B4 exists precisely so
that this trade can be measured rather than asserted.

\textbf{Weight-triggered selection is prior art; retrospective selection is the variation.} Triggering
detection on a stabilised weight change is already done~\cite{R01}. What differs here is that the
trigger does not start a capture --- it selects, from history, a window that \emph{already contains}
the event, along with a baseline from before it, and refuses when that window is not complete.
Section~\ref{sec:timing} shows that this variation is what makes the configuration auditable: because
the window is explicit and the manifest records what was realised, the mismatch between the settlement
timestamp and mechanical impact is \emph{visible in the data} instead of silently shaping the evidence.

\textbf{The settlement-timing result is the most consequential finding of this paper}, and it is
analytical rather than empirical. With the shipped firmware the system cannot observe the descent; it
observes the settled product, over near-duplicate frames, with roughly $0.30$--$0.40$~s of decay
budget before the baseline is contaminated and the system emits a mislabelled false reject. Any
deployment of this configuration must measure $\Delta_s$ first. A reader who takes one engineering
lesson from this manuscript should take this one: in an event-triggered perception system, the
\emph{semantics} of the trigger timestamp --- stable-sample versus impact --- silently determines what
the evidence window contains.

%---------------------------------------------------------------------
\section{Limitations and Threats to Validity}
\label{sec:limits}

\textbf{No empirical validation of the primary claims.} Every accuracy, latency, resource, cost and
loss-reduction figure is pending. Static code inspection is not proof of successful hardware
operation, and a passing unit-test suite over constructed inputs is not evidence about real packaging,
real lighting or real middleware.

\textbf{A narrow visual match is not end-to-end theft prevention.} The visual gate establishes, at
best, that \emph{some} foreground region in a settled-mass window was consistent with the scanned SKU
under a cosine threshold and a competitor margin. It does not establish that no other object entered
the basket, that the matched region is the object that caused the mass change, or that an occluded
exchange did not occur. Scenario~S16 is explicitly out of scope and must remain unresolved rather than
resolved optimistically.

\textbf{The two mass checks are not independent sensors.} They are computed by different processors
from different baselines with different code paths, which protects against a software fault or a stale
expectation on one side. They read the \emph{same load cell through the same HX711}. A mechanical
fault, a load path bypassing the cell, a calibration error or an ADC fault defeats both
simultaneously. Describing this as sensor redundancy would be wrong.

\textbf{Apparent size is not physical size.} Equation~(\ref{eq:aspect}) uses an image-plane aspect
ratio from an uncalibrated monocular camera with no depth channel. It cannot bound physical
dimensions, and no claim depends on it beyond a weak consistency filter.

\textbf{Single aggregate load cell.} Concurrent placements cannot be attributed. The system flags them
as ambiguous, which is correct but is a \emph{capability limit}, not a solution. Multiple cells are
planned, not implemented.

\textbf{Security and payment.} The cart LAN is assumed trusted; DDS security is not enabled, so topic
injection by a host on the same domain is not prevented. A single forged verdict whose $(\tau,a)$ and
barcode match a pending unit would be accepted. Transaction and attempt binding provides \emph{replay}
resistance but not authentication. On payment: the system generates a revision-gated UPI deep-link
QR~\cite{R80}. \textbf{Generating a payment request is not confirmed settlement.} A displayed static
QR cannot revoke a screenshot already taken and cannot prove that money moved. Unattended commercial
checkout additionally requires provider-side order authorisation and a confirmed-settlement webhook,
neither of which exists in this code.

\textbf{Thresholds are unvalidated defaults.} The nine decision thresholds in Table~\ref{tab:params}
were chosen during development. None is derived from an operating-characteristic analysis on held-out
data.

\textbf{Catalogue and generalisation scope.} The schema requires a single fixed gross mass per SKU,
excluding loose produce and variable-weight goods, and the mass channel cannot discriminate items
below the noise gate. Reference photographs must be product crops in the deployment camera's own
domain; references captured on a different camera measure the domain gap rather than the method. A
database built with one embedding backend is meaningless to another, which the code enforces by
refusing a mismatched backend stamp.

\textbf{Internal validity.} The settlement-timing analysis assumes the firmware constants of
Section~\ref{sec:impl}, ${\sim}100$~ms HX711 sample gating, and serial latency small relative to
$0.5$~s. The first two are read from source; the third is an assumption pending measurement. The
bandwidth and memory figures are arithmetic from the configuration and exclude all protocol overhead.

\textbf{External validity.} A single prototype cart, one camera geometry, one mounting stiffness, one
LAN and one small catalogue cannot support store-scale claims. Catalogue growth in particular is
expected to degrade the all-SKU ranking~\cite{R07}, so any reported operating point is specific to the
catalogue size at which it was measured.

\textbf{Construct validity.} The metrics measure billing correctness and gate behaviour. They do not
measure shrinkage reduction, which depends on shopper behaviour, store process and deterrence, and
which cannot be inferred from a bench matrix. No claim about loss reduction is made.

\textbf{Literature-review completeness.} The review draws on a curated pack of 85 resources, of which
10 were read in full. It is \emph{not} a systematic review: no protocol-registered database query,
date range, screening log or citation saturation was performed. Items reviewed only at metadata or
abstract level are labelled in Table~\ref{tab:related}, and no claim depends on an unread full text.
Where supplied papers contained internal inconsistencies these were recorded and their citation
mappings were not reused.

%---------------------------------------------------------------------
\section{Conclusion and Future Work}

We have specified NXTCart, a barcode-first smart shopping cart in which a scanned product's identity
is checked by two independent, conjunctive gates: an independently recomputed settled-mass comparison
on an Arduino/HX711, and a foreground-only visual verification on a separate host over ROS~2. We do
not claim the barcode/weight/image combination, weight reconciliation, or weight-triggered inference
as novel; close prior art exists for each. The design choices we document and defend are mechanical:
retrospective selection of a dense, gap-checked evidence window from a bounded source-timestamped
history around a settled-mass event, together with a pre-entry baseline; embedding only connected
foreground regions so that no scene background ever reaches the network; scoring every region against
every registered SKU so that a confident alternate product vetoes a matching region elsewhere in the
same burst; binding every verdict to a transaction and attempt; and abstaining when the evidence
window is incomplete.

The paper's substantive analytical result is that these mechanisms are only as good as the semantics of
their trigger. Because the settlement timestamp marks a six-sample stability decision rather than a
mechanical impact, the shipped configuration cannot observe the product in descent and is
self-consistent only while the impact-to-settlement delay remains between $0.40$~s and $0.90$~s,
leaving roughly $0.30$--$0.40$~s of mechanical decay budget before the selected baseline is
contaminated by the product itself and the system emits a mislabelled false reject. Measuring that
delay distribution is the first required step of any deployment, and the diagnostic manifests make it
directly checkable.

We also report, as a verified fact, that the software checks pass at the stated revision --- 103
integration and regression tests, interface parity between hosts, and an enforced deployment boundary
--- and we state plainly that this establishes none of the system's accuracy, timing or security
properties.

Future work falls into three tiers, deliberately separated by evidential status.
\emph{Required before any accuracy claim:} execute the protocol of Section~\ref{sec:eval}; measure
$\Delta_s$ with an independent impact reference and re-tune $\delta_{\mathrm{pre}}$,
$\delta_{\mathrm{base}}$ and $T_{\mathrm{buf}}$ together, or publish a distinct impact timestamp;
measure achieved frame rate, gap distribution and wire bandwidth; replace the fixed thresholds with a
reported risk--coverage curve on held-out confusable pairs; complete native \texttt{colcon} and
Arduino builds and live DDS acceptance. \emph{Natural extensions of the implemented design:}
calibrated confidence in place of raw cosine scores~\cite{R60,R61}; stronger frozen backbones or
vision-language retrieval as drop-in replacements~\cite{R58,R59,R62}; a manifest-driven regression
corpus so threshold changes can be replayed offline; drift compensation in the mass channel along the
lines of~\cite{R02}; and provider-side payment authorisation with confirmed settlement~\cite{R80}.
\emph{Proposed but unimplemented, and not to be described otherwise:} monitored-aperture light
curtains, multiple load cells for position-resolved mass, IMU compensation for cart motion, continuous
basket reconciliation, in-store navigation and inventory forecasting~\cite{R21}. These appear in the
project's design material as planned work and have no implementation evidence in the code examined
here.

%---------------------------------------------------------------------
\section*{Acknowledgment}
[ACKNOWLEDGMENT PLACEHOLDER --- funding sources, grant numbers, institutional support, equipment loans
and non-author contributors. The authors thank [NAMES] for [CONTRIBUTION]. Remove this section if not
applicable; do not list any funder that has not approved acknowledgment.]

%---------------------------------------------------------------------
\begin{thebibliography}{99}
\footnotesize

\bibitem{R49} B.~Santra and D.~P.~Mukherjee, ``A comprehensive survey on computer vision based
approaches for automatic identification of products in retail store,'' \emph{Image and Vision
Computing}, vol.~86, pp.~45--63, 2019.

\bibitem{R50} Y.~Wei, S.~Tran, S.~Xu, B.~Kang, and M.~Springer, ``Deep learning for retail product
recognition: Challenges and techniques,'' \emph{Computational Intelligence and Neuroscience},
vol.~2020, pp.~1--23, 2020.

\bibitem{R51} R.~Oucheikh, T.~Pettersson, and T.~L\"ofstr\"om, ``Product verification using OCR
classification and Mondrian conformal prediction,'' \emph{Expert Systems with Applications}, vol.~188,
p.~115942, 2022.

\bibitem{R81} Instacart, ``Caper Carts: Official product overview.'' [Online]. Available:
\url{https://company.instacart.com/enterprise-platform/connected-stores/caper-carts} (accessed
Oct.~12, 2026). Vendor product claim; not peer-reviewed evidence.

\bibitem{R82} Amazon, ``5 ways the newest Dash Cart can save time and money,'' 2026. [Online].
Available: \url{https://www.aboutamazon.com/news/retail/amazon-redesigned-dash-cart-whole-foods-market}
(accessed Oct.~12, 2026). Vendor product claim; not peer-reviewed evidence.

\bibitem{R84} V.~N.~Prithvish, S.~Agrawal, and J.~S.~R.~Alex, ``An IoT-based smart shopping cart using
the contemporary barcode scanner,'' in \emph{Intelligent Embedded Systems}, Lecture Notes in
Electrical Engineering 492, 2018, pp.~45--58.

\bibitem{R85} U.~Gangwal, S.~Roy, and J.~Bapat, ``Smart shopping cart for automated billing purpose
using wireless sensor networks,'' in \emph{SENSORCOMM 2013}, IARIA, 2013, pp.~168--172.

\bibitem{R02} M.~Desai, B.~N.~K., M.~Shirol, K.~A.~Varun, and C.~J., ``A weight-reconciliation
algorithm for fraud detection in smart shopping carts,'' in \emph{Proc. 4th Int. Conf. Knowledge
Engineering and Communication Systems (ICKECS)}, 2026, pp.~1--4.

\bibitem{R01} J.~Zhang, Z.~Xu, C.~Liang, and Y.~Zhu, ``Intelligent shopping cart product recognition
system based on computer vision and multi-sensor fusion,'' in \emph{Proc. 5th Int. Conf. Communication
Technology and Information Technology (ICCTIT)}, 2025, pp.~249--253.

\bibitem{R07} Z.~Al~Adzani~Hidayat, H.~F.~Ramadhan, P.~Hidayatullah, and N.~Syakrani, ``A few-shot
learning pipeline for retail product recognition using object detection and metric embeddings,'' in
\emph{Proc. 8th Int. Conf. Information and Communications Technology (ICOIACT)}, 2025, pp.~42--47.

\bibitem{R52} B.~Ciapas and P.~Treigys, ``Centre-loss---A preferred class verification approach over
sample-to-sample in self-checkout products datasets,'' \emph{IET Computer Vision}, vol.~18, no.~7,
pp.~1004--1016, 2024.

\bibitem{R31} M.~Sandler, A.~Howard, M.~Zhu, A.~Zhmoginov, and L.-C.~Chen, ``MobileNetV2: Inverted
residuals and linear bottlenecks,'' in \emph{Proc. IEEE/CVF Conf. Computer Vision and Pattern
Recognition (CVPR)}, 2018, pp.~4510--4520.

\bibitem{R54} S.~Macenski, T.~Foote, B.~Gerkey, C.~Lalancette, and W.~Woodall, ``Robot Operating
System 2: Design, architecture, and uses in the wild,'' \emph{Science Robotics}, vol.~7, no.~66, 2022.

\bibitem{R61} Y.~Geifman and R.~El-Yaniv, ``Selective classification for deep neural networks,'' in
\emph{Advances in Neural Information Processing Systems 30}, 2017.

\bibitem{R10} M.~Shahroz, M.~F.~Mushtaq, M.~Ahmad, S.~Ullah, A.~Mehmood, and G.~S.~Choi, ``IoT-based
smart shopping cart using radio frequency identification,'' \emph{IEEE Access}, vol.~8,
pp.~68426--68438, 2020.

\bibitem{R18} P.~Chandrasekar and T.~Sangeetha, ``Smart shopping cart with automatic billing system
through RFID and ZigBee,'' in \emph{Proc. Int. Conf. Information Communication and Embedded Systems
(ICICES)}, 2014, pp.~1--4.

\bibitem{R11} P.~A.~Jeyanthy, S.~R.~Mohamed, M.~Alarfaj, S.~S.~Reddy, and K.~S.~Jenin, ``Design and
development of a smart shopping cart system,'' in \emph{Proc. 3rd Int. Conf. Intelligent Computing
Instrumentation and Control Technologies (ICICICT)}, 2022, pp.~184--188.

\bibitem{R12} S.~Sakthivel, A.~R.~R.~K., S.~Sai~B., V.~K.~K.~M., and G.~Kumar~C., ``AIoT-based smart
cart system,'' in \emph{Proc. Int. Conf. Smart Technologies and Systems for Next Generation Computing
(ICSTSN)}, 2022, pp.~1--6.

\bibitem{R83} S.~Karjol, A.~K.~Holla, C.~B.~Abhilash, P.~V.~Amrutha, and Y.~V.~Manohar, ``An IoT based
smart shopping cart for smart shopping,'' in \emph{Cognitive Computing and Information Processing},
Communications in Computer and Information Science 801, 2018, pp.~373--385. Cited with the publisher's
corrected five-author attribution.

\bibitem{R06} R.~Li, T.~Song, N.~Capurso, J.~Yu, J.~Couture, and X.~Cheng, ``IoT applications on
secure smart shopping system,'' \emph{IEEE Internet of Things Journal}, vol.~4, no.~6,
pp.~1945--1954, 2017.

\bibitem{R47} B.~K.~Rajithkumar, G.~M.~Deepak, B.~V.~Uma, B.~N.~Hadimani, A.~R.~Darshan, and
C.~R.~Kamble, ``Design and development of weight sensors based smart shopping cart and rack system for
shopping malls,'' \emph{Materials Today: Proceedings}, vol.~5, no.~4, pp.~10814--10820, 2018.

\bibitem{R48} A.~Yewatkar, F.~Inamdar, R.~Singh, Ayushya, and A.~Bandal, ``Smart cart with automatic
billing, product information, product recommendation using RFID and ZigBee with anti-theft,''
\emph{Procedia Computer Science}, vol.~79, pp.~793--800, 2016.

\bibitem{R17} B.~Chaure and P.~Jain, ``Development of e-shopping cart with theft control mechanism: No
queue,'' in \emph{Proc. Int. Conf. Emerging Technological Trends (ICETT)}, 2016, pp.~1--5.

\bibitem{R04} N.~Ragesh, B.~Giridhar, D.~Lingeshwaran, P.~Siddharth, and K.~P.~Peeyush, ``Deep
learning based automated billing cart,'' in \emph{Proc. Int. Conf. Communication and Signal Processing
(ICCSP)}, 2019, pp.~779--782.

\bibitem{R05} T.~W.~Alessa, F.~Ahmed~Alfarhan, F.~K.~Bugis, A.~Abdulhadi~Aldawood, S.~A.~Alhuwaidi,
and A.~Abul~Hussain, ``Smart shopping cart with machine learning, navigation, and weighing
technologies,'' in \emph{Proc. 6th Int. Symp. Advanced Electrical and Communication Technologies
(ISAECT)}, 2024, pp.~1--6.

\bibitem{R03} A.~Naithani, T.~S.~Rajpal, and V.~Mistry, ``TAV-Cart: Transaction assistance and
validation cart system for effortless shopping and checkout,'' in \emph{Proc. Int. Conf. Advances in
Computing, Communication and Materials (ICACCM)}, 2024, pp.~1--6.

\bibitem{R26} A.~Prabhath, A.~Anoop, and V.~Sureshkumar, ``Multi-modal product recognition and
in-store navigation using an intelligent shopping cart,'' in \emph{Proc. 5th Int. Conf. Computer
Networks, Big Data and IoT (ICCBI)}, 2026, pp.~1241--1245.

\bibitem{R09} M.~A.~Sarwar, Y.-A.~Daraghmi, K.-W.~Liu, H.-C.~Chi, T.-U.~Ik, and Y.-L.~Li, ``Smart
shopping carts based on mobile computing and deep learning cloud services,'' in \emph{Proc. IEEE
Wireless Communications and Networking Conf. (WCNC)}, 2020, pp.~1--6.

\bibitem{R24} Y.~Berdaliyev and A.~P.~James, ``RFID-cloud smart cart system,'' in \emph{Proc. Int.
Conf. Advances in Computing, Communications and Informatics (ICACCI)}, 2016, pp.~2346--2352.

\bibitem{R21} M.~I.~Zulfiqar, A.~Khalid, A.~Siddig, M.~J.~Nawaz, and S.~Saay, ``AI-driven smart
shopping carts with real-time tracking and inventory forecasting for enhanced retail efficiency,''
\emph{IEEE Access}, vol.~13, pp.~55576--55585, 2025.

\bibitem{R27} L.~H.~Pham, D.~Nguyen-Ngoc~Tran, H.-H.~Nguyen, T.~H.-P.~Tran, H.-J.~Jeon, H.-M.~Jeon,
and J.~W.~Jeon, ``DeepACO: A robust deep learning-based automatic checkout system,'' in \emph{Proc.
IEEE/CVF Conf. Computer Vision and Pattern Recognition Workshops (CVPRW)}, 2022, pp.~3106--3113.

\bibitem{R42} X.-S.~Wei, Q.~Cui, L.~Yang, P.~Wang, L.~Liu, and J.~Yang, ``RPC: A large-scale and
fine-grained retail product checkout dataset,'' \emph{Science China Information Sciences}, vol.~65,
no.~9, 2022.

\bibitem{R38} E.~Goldman, R.~Herzig, A.~Eisenschtat, J.~Goldberger, and T.~Hassner, ``Precise
detection in densely packed scenes,'' in \emph{Proc. IEEE/CVF Conf. Computer Vision and Pattern
Recognition (CVPR)}, 2019, pp.~5222--5231.

\bibitem{R39} M.~Klasson, C.~Zhang, and H.~Kjellstr\"om, ``A hierarchical grocery store image dataset
with visual and semantic labels,'' in \emph{Proc. IEEE Winter Conf. Applications of Computer Vision
(WACV)}, 2019, pp.~491--500.

\bibitem{R63} Y.~Yao, X.~Tian, Z.~Tang, S.~Biswas, H.~Lei, T.~Gedeon, and L.~Zheng, ``Training with
product digital twins for AutoRetail checkout,'' arXiv:2308.09708, 2023.

\bibitem{R08} J.-Y.~Jeon, S.-W.~Kang, H.-J.~Lee, and J.-S.~Kim, ``A retail object classification
method using multiple cameras for vision-based unmanned kiosks,'' \emph{IEEE Sensors Journal},
vol.~22, no.~22, pp.~22200--22209, 2022.

\bibitem{R30} I.~Baz, E.~Yoruk, and M.~Cetin, ``Context-aware hybrid classification system for
fine-grained retail product recognition,'' in \emph{Proc. IEEE 12th Image, Video, and
Multidimensional Signal Processing Workshop (IVMSP)}, 2016, pp.~1--5.

\bibitem{R58} A.~Radford \emph{et al.}, ``Learning transferable visual models from natural language
supervision,'' in \emph{Proc. 38th Int. Conf. Machine Learning (ICML)}, 2021.

\bibitem{R59} M.~Oquab \emph{et al.}, ``DINOv2: Learning robust visual features without supervision,''
arXiv:2304.07193, 2023.

\bibitem{R62} E.~G.~Maminta and R.~O.~Atienza, ``What matters for grocery product retrieval with open
source vision language models,'' arXiv:2605.18029, 2026. Cited as the verified arXiv preprint version;
no final proceedings record is claimed.

\bibitem{R45} J.~Falc\~{a}o, C.~Ruiz, S.~Pan, H.~Y.~Noh, and P.~Zhang, ``FAIM: Vision and weight
sensing fusion framework for autonomous inventory monitoring in convenience stores,'' \emph{Frontiers
in Built Environment}, vol.~6, 2020.

\bibitem{R46} C.~Ruiz, J.~Falcao, S.~Pan, H.~Y.~Noh, and P.~Zhang, ``AIM3S: Autonomous inventory
monitoring through multi-modal sensing for cashier-less convenience stores,'' in \emph{Proc. 6th ACM
Int. Conf. Systems for Energy-Efficient Buildings, Cities, and Transportation}, 2019, pp.~135--144.

\bibitem{R37} D.~L.~Hall and J.~Llinas, ``An introduction to multisensor data fusion,''
\emph{Proceedings of the IEEE}, vol.~85, no.~1, pp.~6--23, 1997.

\bibitem{R53} B.~Khaleghi, A.~Khamis, F.~O.~Karray, and S.~N.~Razavi, ``Multisensor data fusion: A
review of the state-of-the-art,'' \emph{Information Fusion}, vol.~14, no.~1, pp.~28--44, 2013.

\bibitem{R60} C.~Guo, G.~Pleiss, Y.~Sun, and K.~Q.~Weinberger, ``On calibration of modern neural
networks,'' in \emph{Proc. 34th Int. Conf. Machine Learning (ICML)}, 2017.

\bibitem{R34} J.~Redmon, S.~Divvala, R.~Girshick, and A.~Farhadi, ``You only look once: Unified,
real-time object detection,'' in \emph{Proc. IEEE Conf. Computer Vision and Pattern Recognition
(CVPR)}, 2016, pp.~779--788.

\bibitem{R35} A.~Bewley, Z.~Ge, L.~Ott, F.~Ramos, and B.~Upcroft, ``Simple online and realtime
tracking,'' in \emph{Proc. IEEE Int. Conf. Image Processing (ICIP)}, 2016, pp.~3464--3468.

\bibitem{R36} N.~Wojke, A.~Bewley, and D.~Paulus, ``Simple online and realtime tracking with a deep
association metric,'' in \emph{Proc. IEEE Int. Conf. Image Processing (ICIP)}, 2017, pp.~3645--3649.

\bibitem{R72} Open Robotics, ``ROS~2 Jazzy: About quality of service settings.'' [Online]. Available:
\url{https://docs.ros.org/en/jazzy/Concepts/Intermediate/About-Quality-of-Service-Settings.html}
(accessed Oct.~12, 2026).

\bibitem{R74} Avia Semiconductor, ``HX711: 24-bit analog-to-digital converter for weigh scales,''
datasheet. [Online]. Available: \url{https://cdn.sparkfun.com/assets/b/f/5/a/e/hx711F_EN.pdf}
(accessed Oct.~12, 2026).

\bibitem{R75} B.~Necula and contributors, ``HX711 Arduino library.'' [Online]. Available:
\url{https://github.com/bogde/HX711} (accessed Oct.~12, 2026).

\bibitem{R79} Arduino, ``A guide to EEPROM.'' [Online]. Available:
\url{https://github.com/arduino/docs-content/blob/main/content/learn/03.programming/07.eeprom-guide/eeprom-guide.md}
(accessed Oct.~12, 2026).

\bibitem{R78} Natural History Museum and contributors, ``pyzbar: Barcode and QR decoding.'' [Online].
Available: \url{https://github.com/NaturalHistoryMuseum/pyzbar} (accessed Oct.~12, 2026).

\bibitem{R77} OpenCV, ``Deep learning in OpenCV.'' [Online]. Available:
\url{https://github.com/opencv/opencv/wiki/Deep-Learning-in-OpenCV} (accessed Oct.~12, 2026).

\bibitem{R66} R.~Szeliski, \emph{Computer Vision: Algorithms and Applications}, 2nd~ed. Springer,
2022.

\bibitem{R80} National Payments Corporation of India, ``UPI frequently asked questions.'' [Online].
Available: \url{https://www.npci.org.in/what-we-do/upi/faqs} (accessed Oct.~12, 2026).

\end{thebibliography}

\end{document}
